\section{Canonical names and conditional observations}
\label{names:section}

This section constructs the information variables used in the positive-cost
cylinder argument. There are two distinct uses of probability. Canonical
entropies concern the assigned law of a single experiment. Independent
observations and the walks constructed below are auxiliary experiments,
conditioned on a specified common side variable. We shall keep their laws
separate.

Throughout, we work on an affine extremal configuration satisfying
\eqref{eq:F}, with ordinary regularity imposed before the documented gates.
Thus
\[
 r_1=f_1=1,\qquad p+d=1+3d-2\Gamma.
\]
In the hybrid case the growth exponent is strictly larger than the
classical comparison exponent, so that \eqref{eq:Pl} is available.
The limiting conventions of the framework apply: every depth gap, finite
query list, and operator stack is fixed before the preceding approximation
errors tend to zero. In particular, an error denoted by \(o(h)\) at a
tangent gap \(h\) is obtained by taking those earlier limits first.

\subsection{The canonical coordinates}

For a real number or vector \(z\), let \([z]_b\) be its absolute dyadic
cell at width comparable to \(s^b\), or the fixed representative of that
cell when a coordinate is required. The convention applies also to
negative \(b\). We also write \(X^{(v)}\) for the angular cell
\([X]_v\). Let \(t_a\) be the left endpoint of the depth-\(a\)
ancestor of the terminal time. Approximate gates may be used, with their
errors taken to zero before any subsequent gap is reduced.

In the classical model, \(y_0\) denotes the exact particle, including all
exact auxiliary indices. A positive-time label on its realized path is
determined by \(y_0\) and the corresponding interval start. The label at
the final assignment is still sampled. In the hybrid model the base
coordinates are exact; the sampled normal coordinates are extended
backward as a classical ray only for the purpose of defining names.
Comparisons of this artificial ray with physical rows will always use
\eqref{eq:Pl}.

A canonical name is a finite record of a ray in a specified shear frame.
It stores the rounded coefficients of that frame, then the base coordinates
and the residual normal coordinates. For example, a linear shear with
recorded coefficient $L$ replaces normal velocity and position by
$N-LX$ and $D_{t_a}-LB_{t_a}$. Their bins describe a cell in that known
frame; they do not specify an exact particle. The displays below choose
which curvature coefficients to retain, and the templates choose the
precision of the entries.

Finite partitions and their conditional entropies across resolutions also
underlie the local entropy methods of Hochman--Shmerkin
\cite[Section~4]{HochmanShmerkin2012}. Here the records must additionally retain
the time gate and the changing shear.

\begin{definition}[Curvature displays]\label{names:curvature-displays}
Fix one of the following displays on an entire parameter patch.
\begin{enumerate}[label=(\roman*)]
\item In the quadratic equality case, write
\[
 b(a)=b_2(a)=c(a-1),\qquad b_1(a)=(c-J)(a-1),\qquad c>0,
\]
and put \(P^\diamond=P_1(F)\). If \(c>1\), let \(R_a,k_a\)
be the rounded values of \(\beta_1(t_a),K_1\), at depths
\(b(a),b(a)-a\), respectively, and extend \(R_a\) with slope
\(k_a\) from \(t_a\). If \(c\leq1\), round only
\(\beta^\diamond=\beta_1(T_1)\) at depth \(b(a)\), and set
\(k_a=0\).
\item In the pure linear case \(c_2=0<c_1\), set \(R_a=k_a=0\),
\(b_1(a)=c_1(a-1)\), and take \(P^\diamond\) to be the
linearized terminal coefficient at the center of its own base test.
Earlier tests are replaced by their enclosing linear tests from
\eqref{eq:F}, with vanishing exponent loss.
\item In the strict case \(c_2=1\), \(q<1\), fix an endpoint
\(D\in(0,1]\); hybrid applications use \(D=1\). Use the encountered
label \(\lambda_D\), put \(P^\diamond=P_D(F)\), and first record
\[
 \bar k=[K_D]_{b_2(D)-D}.
\]
Subtract the common shear whose curvature trajectory is
\((\bar k t,\bar k)\), with zero other coefficients. All subsequent
coordinates in this display are the residual coordinates. Set
\[
 b(a)=b_2(D)+a-D\quad(a\leq D),
\]
round the residual \(\beta_D(T_D)\) at depth \(b(a)\) to obtain
\(R_a\), and set \(k_a=0\). The actual derivative exponent remains
\(b_1(a)=(q-J)(a-1)\). For \(D<1\) the law remains the final
classical assigned law; its encountered \(D\)-label supplies the data.
\end{enumerate}
Write \(P\) for the specified derivative, after the global subtraction
in case (iii).
\end{definition}

The horizontal resolution is \(v\geq0\). We use one of the templates
\[
\begin{array}{c|c|c}
 \text{template}&l(a,v)&w(a,v)\\ \hline
 \text{quadratic}&b(a)+v&b(a)+2v\\
 \text{clipped}&\min\{b(a)+v,b_1(a)\}&l(a,v)+v\\
 \text{pure linear}&b_1(a)&b_1(a)+v.
\end{array}
\]
The clipped template is used only when \(b_1'>0\). All derivatives of
\(b\), or of \(b_1\) when it occurs as an active branch, are positive.
An explicitly indicated enlargement parameter \(g\geq0\) changes only
the normal phase widths from \(w,a+w\) to \(w-g,a+w-g\).
Unless indicated otherwise, \(g=0\).

\begin{definition}[Canonical name]\label{names:canonical-name}
The name \(\mathcal C(a,v)\) records \(\bar k\) when present,
\(t_a,R_a,k_a\), and
\[
 L=[P-2R_aX+2k_aB_{t_a}]_l.
\]
It also records
\[
 [X]_v,\qquad [B_{t_a}]_{a+v},\qquad
 [Y_{t_a}-2[X]_vB_{t_a}]_{a+2v}.
\]
Finally it records the residual normal velocity and position at \(t_a\),
at depths \(w-g,a+w-g\), after subtracting the shear
\[
\begin{split}
 N&\longmapsto N-LX-R_aX^2-k_a(Y_{t_a}-2XB_{t_a}),\\
 D_{t_a}&\longmapsto D_{t_a}-LB_{t_a}-R_aY_{t_a}+k_aB_{t_a}^2.
\end{split}
\]
Here \(R_a\) in the displayed expressions is its value at \(t_a\).
\end{definition}

All names are used on significant polynomial support. The cutoff
convention permits us to remove negligible mass and enlarge a fixed
polynomial support on each inner setup. Thus finite names have the
cardinality bounds needed for likelihood and entropy arguments. Exact
side variables need not have finite range. Conditional kernels can be
taken on the standard Borel codes of
Lemma~\ref{cyl:borel-paths}. Completing-path tests below are
projections of the specified Borel finite-path relations. Before
another path test uses a completed-measurable row restriction, we
take its Borel version under that stage's active row law as in the
lemma.

\begin{lemma}[Refinement and comparison]\label{names:refinement}
Fix a template and \(g\). If both \(a\) and \(v\) increase, the finer
canonical name determines the coarser one up to bounded lists. The
relative branching exponent is bounded by a constant times the sum of
the two depth changes. At approximated gates the lists may have size
\(s^{-o(1)}\). Names at nearby parameters admit mutual comparison
lists with exponent tending to zero with the parameter displacement.
\end{lemma}

\begin{proof}
The time cells are nested. For base coordinates, backflow uses
\(B_t=B_{t'}+(t-t')X\), \(Y_t=Y_{t'}+(t-t')X^2\).
After subtracting the tilt \(2[X]_vB_t\), the unresolved contribution
of the time displacement is quadratic in the unresolved angular
increment. Its sizes are therefore \(s^{a+v}\) and \(s^{a+2v}\).
Changing the angular cell center produces its known linear correction
and an error of the latter size. These facts give both the base lists
and their relative branching bounds.

Compare two curvature displays in the coarse time units. Their
differences have sizes \(O(s^{b(a)})\) and
\(O(s^{b(a)-a})\). For a nonconstant display, extrapolation backward
from the finer time preserves these bounds because \(b'-1\geq0\).
Constant displays require no slope comparison. After both displays
have been guessed, the unresolved part of their derivative difference
has size at most \(s^{b(a)+v}\), which is no larger than \(s^l\).
Thus the derivative cells have bounded comparison lists.

For completeness, the normal estimate uses the difference of the two
shears, rather than either shear separately. Its derivative on the
actual base ray is \(O(s^l)\), since both displayed derivatives
approximate \(P\). Its curvature difference has size
\(O(s^{b(a)})\) in the coarse time units. On the unresolved base cell,
Taylor's identity bounds its velocity uncertainty by
\[
 O(s^{l+v}+s^{b(a)+2v})=O(s^w),
\]
and its position uncertainty by \(s^a\) times this quantity.
These estimates are stronger when the normal widths are enlarged.
Shear subtraction commutes with backflow. This proves the lists for
the residual coordinates even when the absolute shear or absolute
position is large. The same calculations with unequal resolutions
give at most \(s^{-C(|\Delta a|+|\Delta v|)-o(1)}\) choices.
Comparison through a common coarser pair proves the final assertion.
\end{proof}

\subsection{Entropy bounds and label prediction}

We request deterministic final-readout profiles for the finite names
used below. Let \(\mathcal H(a,v)\) denote the limiting entropy of
\(\mathcal C(a,v)\), divided by \(\log(1/s)\), and put
\[
 \mathsf D_0=\lim\log_{1/s}(e/\kappa),\qquad
 W(a,v)=(3+\eta)v+(1-\eta)(w(a,v)-g).
\]

\begin{proposition}[Canonical entropy and endpoint recovery]
\label{names:canonical-entropy}
On its domain of use, the excess satisfies
\begin{equation}
 E(a,v)=\mathcal H(a,v)-\mathsf D_0-(p+d)a-W(a,v)\geq0.
 \tag{canon}\label{eq:canon}
\end{equation}
For hybrid names this assertion is used only when \(g=0,w<0\), and
at their terminal limits by continuity. The following equality and
prediction statements also hold, with lists of size \(s^{-o(1)}\).
\begin{enumerate}[label=(\roman*)]
\item For \(D=1,g=0\), equality holds at \((1,0)\). In the
classical equality and pure linear cases with \(b_1'>0\), equality
holds at \(v=u(a)\). In a classical strict family ending at \(D\),
it holds at \((D,u(D))\). At these equality points use the template
with \(l=b+v=b_1\), or the pure linear template.
\item In a strict family every row connected from \(i<D\) to \(D\),
and its label, predicts \(\bar k\). After its subtraction,
\[
 |K_i|\leq s^{b_2(D)-D-o(1)},\qquad
 |\beta_D(T_D)-\beta_i(T_i)|\leq s^{b(i)-o(1)}.
\]
\item In an equality or pure linear family, suppose
\(J>0\), \(a<j<1\), \(v=u(j)>0\), and
\[
 l(a,v)\leq\min_{[j,1]}b_1.
\]
Then \(\lambda_j\) predicts \(\mathcal C(a,v)\) on every completing
path; in the hybrid case require \(w<0\).
\item In the classical cases with \(b_1'>0\), fix \(0<a<j\).
If \(v-u(j)\geq0\) is sufficiently small compared with \(j-a\)
and \(l(a,v)\leq b_1(j)\), then
\((\lambda_j,X^{(v)})\) predicts \(\mathcal C(a,v)\).
The strict version uses \(j=D\). There \(\lambda_D\) alone suffices
at \(v=u(D)\); this last assertion also holds for hybrid \(D=1\)
and \(w<0\), without the condition \(b_1'>0\).
\end{enumerate}
No equality is asserted at an arbitrary point where
\(\min_{[a,1]}b_1<b_1(a)\).
\end{proposition}

\begin{proof}
Fix a canonical name at depth \(a\). On its ancestor rows the base
coordinates lie in the displayed base cell; after the displayed shear
the two normal coordinates lie in cells of widths
\(s^{w-g},s^{a+w-g}\). Conversion between \(t_a\) and the actual
gate time has the same precision because the residual velocity is
known to that precision. Thus those rows are covered by a small list
of regularity tests of weight \(s^W\), in a small list of time
intervals. Their benchmark mass is at most
\(s^{pa+W-o(1)}\kappa\). Division by the assigned-law benchmark
\(s^{-da}e\), and the sparse-list consequence of \eqref{eq:Sp},
exclude cells whose typical probability is larger than
\((\kappa/e)s^{(p+d)a+W-o(1)}\). This is \eqref{eq:canon}.
In the hybrid comparison the artificial ray differs from the ancestor
row by at most \(s^{-o(1)}\) in normal velocity and
\(s^{a-o(1)}\) in position, by \eqref{eq:Pl}. Since the base is
identical, all shears preserve this additive comparison. It fits the
required widths when \(w<0\). The classical virtual endpoint uses
the same sparse-list argument on the final law.

For the upper bounds at the stated equality points, the reached
label determines the name to a small list. Curvature comparisons are
those of \eqref{eq:F}. If \(c\leq1\), omitting the slope of the
curvature trajectory costs at most the own curvature width in its
own duration units. The derivative discrepancy along an actual ray
is bounded by \(s^{\min b_1-o(1)}\) on the connecting interval.
Once the curvature cuts are guessed, recentering over the label's
base cell adds at most \(s^{b_2(a)+u(a)-o(1)}\) to that derivative.
The Taylor calculation in Lemma~\ref{names:refinement} then gives the
normal widths. In a strict family the residual own slope is bounded
by \(s^{b_2(D)-D-o(1)}\), so its omission is harmless at \(D\).
The documented label count supplies the reverse entropy inequality.

Assertion (ii) follows directly from \eqref{eq:A}. The slope
comparison on \([i,D]\) predicts the global bin and gives the first
bound. Multiplying this residual slope bound by the earlier duration
gives the artificial depth \(b(i)\). The coefficient comparison
itself is finer, since \(\min_{[i,D]}b_2>b(i)\), and gives the
second bound.

For (iii) and (iv), first guess the earlier curvature cuts. Backflow
of the later base test, with \(X^{(v)}\) when required, gives the
earlier base widths: its uncertainties \(s^{j+u(j)}\) and
\(s^{j+2u(j)}\) fit when the right excess \(v-u(j)\) is small
compared with \(j-a\). At zero excess the inequalities are weak
and the allowed small lists suffice. At the later time the
unresolved angular part of the relative derivative has size at most
\(s^{b(a)+v-o(1)}\); the unresolved position part of a curvature
slope has the additional later-time factor. The remaining derivative
discrepancy fits by the stated condition on \(b_1\). Taylor's
identity now gives the normal comparison, using
\(n(j)\geq w(a,v)\), with strict slack for a small positive excess.

When a curvature slope was omitted, perform this comparison at
\(T_j\), not across the whole earlier duration. The omitted slope
has size \(s^{b_2(j)-j-o(1)}\); its contribution from unresolved
own positions to the derivative is
\(s^{b_2(j)+u(j)-o(1)}\). It fits the later normal widths.
The unresolved angular increment multiplies only the curvature
\emph{difference}, of size \(s^{b(a)-o(1)}\). Having compared
velocity and position at \(T_j\), backflow gives the earlier
position assertion. These bounds depend only on the later label
and the allowed common base region, so the lists are uniform over
all completing paths of that label.
\end{proof}

\subsection{Deterministic profiles and their traces}

The final-readout flattening is performed on growing finite lists of
queries, after the ordinary gates have been prepared in forward
order. Requests may include canonical names, angular bins, mixed
time and coefficient bins, and exact side data such as \(y_0\).
Each retained class has exponent-zero fraction by the count
constraint and minimality of \(p\); subsequent final restrictions
preserve the earlier deterministic profiles. A virtual endpoint
chosen after preliminary profiles is fixed before requesting its
additional profiles. Its exact side variable is common over that
request; no interpolation across unrelated virtual labels is used.
The same procedure applies to one intermediate label selected after
the preliminary profiles: only the final readout is refined again.

For any fixed finite query list and any fixed positive tolerance,
the profile assertions hold off power-small sets. Indeed the
flattening fixes the preselection log probabilities; likelihood
comparison bounds the change under an exponent-zero restriction,
and list counting controls exceptionally small probabilities.
Consequently such deterministic profiles persist under a later
restriction of fraction \(s^{o(h)}\), with errors vanishing in
\(h\)-units after the earlier limits. This fact does not assert
stability of an arbitrary mutual-information budget under that
restriction.

Let
\[
 U(x;a,v)=\lim\frac{-\log\Prob(X^{(x)}\mid\mathcal C(a,v))}
                         {\log(1/s)}
\]
denote the deterministic angular score. It also equals the limiting
conditional entropy. It is nonnegative, nondecreasing and
\(1\)-Lipschitz in \(x\), and vanishes at \(x=v\).
Lemma~\ref{names:refinement} gives Lipschitz dependence on \((a,v)\).
Its increments in \(x\) decrease when either conditioning depth
increases.

\begin{lemma}[Angular trace and time stationarity]
\label{names:angular-trace}
At almost every interior \((a,v)\), and at almost every \(v\) for
each fixed admissible \(a\), there is \(m(a,v)\in[0,1]\) such that
\begin{equation}
 U(v+p'h;a,v)=m(a,v)p'h+o(h)\qquad(p'\geq0).
 \tag{ang}\label{eq:ang}
\end{equation}
The trace has a version nonincreasing in \(a\) at fixed \(v\).
Moreover, at almost every interior root, and at almost every time on
an admissible fixed-\(v\) boundary interval,
\begin{equation}
 U(v+Gh;a,v)-U(v+Gh;a+Gh,v)=o(h)
 \tag{ng}\label{eq:ng}
\end{equation}
for every fixed \(G\), as \(h\downarrow0\).
\end{lemma}

\begin{proof}
At a fixed \(a\), take derivatives in \(x\) for a countable dense
set of conditioning depths. Their values lie in \([0,1]\), and
are nonincreasing as the conditioning depth increases. Monotone
envelopes therefore define the trace from conditioning depths below
\(x\). For \(x\in[v,v+h]\), its derivative at conditioning depth
\(v\) lies between that trace and the derivative at depth
\(x-\eps\), provided \(h<\eps\). At common Lebesgue points,
integration over \([v,v+h]\), followed by \(h\downarrow0\) and
then \(\eps\downarrow0\), makes the two bounds agree. This proves
\eqref{eq:ang}. Countable envelopes and continuity recover the
original profiles. Time monotonicity follows from the corresponding
monotonicity of every difference quotient.

For fixed dyadic \(h\), the function
\(a\mapsto U(v+Gh;a,v)\) is nonincreasing and bounded by \(Gh\).
The integral of its decrement over a translation of length \(Gh\)
is consequently \(O_G(h^2)\) on a compact interval. Divide by
\(h\) and sum over dyadic \(h\). The resulting summable
nonnegative functions tend to zero almost everywhere. Nesting bounds
an arbitrary smaller gap by a comparable dyadic gap. The same proof
can be integrated in \(v\), giving the interior statement.
For a fraction-\(s^{o(h)}\) final restriction, first transfer the
deterministic scores. Chain rule and bounded relative name
capacities then recover the normalized entropy difference; direct
conditioning of the old information budget is not required.
\end{proof}

\begin{lemma}[Mixed coefficient trace]\label{names:mixed-trace}
Let \(V_*\) be an absolute scalar parameter and \(S\) a fixed side
variable. Flatten
\[
 H(x,z)=\lim H_s([V_*]_x\mid S,t_z).
\]
The \(x\)-derivatives lie in \([0,1]\) and decrease with \(z\).
On an affine graph \(x=B(z)\) of slope \(c'>0\), there are
almost-everywhere traces \(0\leq\nu_-(z)\leq\nu(z)\leq1\).
More precisely, at almost every \(a\), fix offsets
\(\alpha_1\leq\alpha_2\) and \(\beta\) for which the following
queries are admissible, and put \(I=[\alpha_1,\alpha_2]\). Then
\[
\begin{split}
 &H(B(a)+\alpha_2h,a+\beta h)
       -H(B(a)+\alpha_1h,a+\beta h)\\
 &\quad=h\bigl(\nu_-(a)|I\cap(-\infty,c'\beta)|
             +\nu(a)|I\cap(c'\beta,\infty)|\bigr)+o(h).
\end{split}
\]
The coefficient-depth interval is split at \(x=B(a+\beta h)\);
the displayed lengths are in units of \(h\). The larger trace is taken from
earlier conditioning times at a fixed coefficient depth.
\end{lemma}

\begin{proof}
Use monotone envelopes of the derivatives for countably many time
depths. At a point of the graph, squeeze the integrand on each
side between its trace and its value with time moved by a fixed
\(\eps\) to that side. Lebesgue differentiation in \(x\), the
positive affine change of variable \(x=B(z)\), and then
\(\eps\downarrow0\) prove the assertion. Joint and conditional
scores agree with these deterministic entropy increments by chain
rule and flattening. This argument uses no regularity of a
coefficient as a function of the exact particle.
\end{proof}

For later classical applications define \(p_*(a)\) to be the
supremum of depths with zero limiting conditional score for the
\emph{raw final} derivative \(P_1(F)\), given \(y_0,t_a\).
Use the raw final linear coefficient in the pure linear case.
Then \(p_*\) is nondecreasing; if \(b_1'>0\),
\(p_*(a)\geq b_1(a)\), since the earlier owned label and
\eqref{eq:A} predict that accuracy. Below \(p_*\), typical values
admit lists of arbitrarily small exponent; strictly above it their
deterministic score is positive. At a virtual endpoint with
\(b_1'>0\), raw \(P_D(F)\) and \(P_1(F)\) differ by at most
\(s^{b_1(D)-o(1)}\), so these conclusions transfer at strictly
coarser query depths.

\begin{proposition}[Memory consequences]\label{names:memory}
Put \(k_0=1-d\) and
\[
 L_{\rm raw}=\max\{1,(1+Q)/2,1+C_{\rm rate}\},\qquad
 (Q,C_{\rm rate})=
 \begin{cases}(q,q-J),&\text{curvature families},\\
 (0,c_1),&\text{pure linear}.
 \end{cases}
\]
Whenever the left label comparison in
Proposition~\ref{names:canonical-entropy}(iii) applies at
\(v=u(j)\), the angular trace satisfies
\[
 m(a,v)\geq
 \frac{k_0+2\Gamma+\eta(L_{\rm raw}-1)}{L_{\rm raw}}.
\]
Fix a classical depth \(j\), and let \(\lambda_j\) denote the
random label encountered there on the final assigned law. Define
\(R_j(v)=\lim H_s(X^{(v)}\mid\lambda_j)\), with conditional
entropy averaged over that label random variable.
Then \(R_j(u(j))=0\), and
\begin{equation}
 \liminf_{z\downarrow0}\frac{R_j(u(j)+z)}z
 \geq\frac{k_0+2\Gamma+\eta(L_{\rm raw}-1)}{L_{\rm raw}}.
 \tag{mr}\label{eq:mr}
\end{equation}
These lower bounds hold on the same flattened law at every tested
boundary prefix whenever \(\lambda_j\) predicts the boundary name;
they are not assertions after conditioning on each individual label
value. Under the small
right-excess comparison, \(m(a,v)\geq R_j'(v)\) almost everywhere.
At typical classical \(j\) with \(C_{\rm rate}>0\), the following
improvements are available:
\begin{enumerate}[label=(\roman*)]
\item If \(p_*(j)>b_1(j)\), with the extra prediction persisting
immediately to the left, replace \(L_{\rm raw}\) by
\(\max\{1,(1+Q)/2\}\).
\item If \(Q\leq1\) and \(0<C_{\rm rate}<1\), an additional
lower bound is \(k_0+2\Gamma-C_{\rm rate}\nu(j)\), where
\(\nu\) is the mixed trace for the raw derivative, side \(y_0\),
and graph \(B=b_1\).
\end{enumerate}
\end{proposition}

\begin{proof}
Pull a proposed angular list back to \(\lambda_j\) using its
canonical prediction list. Apply \eqref{eq:mem} with time gap
\(h=z/L_{\rm raw}\), then transfer the resulting rarity to the
final law by \eqref{eq:Sp}. This proves the prefix statements and
their derivative consequences. Conditioning additionally on
\(\lambda_j,X^{(v)}\) proves the right-excess inequality.

For the prediction variant of \eqref{eq:mem} on \([j-h,j]\),
only a list for the slope at depth \(b_1(j)\) is needed, given
\(y_0,t_{j-h}\). If strict extra prediction persists, this list
has arbitrarily small exponent in \(s^h\). In case (ii), its
cost is at most \(C_{\rm rate}\nu(j)+o(1)\): the score at
\(B(j-h)\) vanishes, and the remaining interval is above the
graph at the earlier conditioning time. Apply
Lemma~\ref{names:mixed-trace}. Restrict the final event to
successful slope predictions, whose failure is negligible at a
fixed strict tolerance. Their lists are then available on every
tested completing path. The prediction memory bound and
\eqref{eq:Sp} prove both enhancements.
\end{proof}

\subsection{Common side variables and orientation marks}

We now arrange for several observations to share a common phase, to the
accuracy of the tangent queries. The side variable below specifies this
common frame; independence will hold after conditioning on that variable.

Write $D_{\rm end}=1$ for the quadratic-equality and pure-linear
displays, and $D_{\rm end}=D$ for the strict-curvature display of
Definition~\ref{names:curvature-displays}.
We now work at \((a,v)\), with \(a>0\) strictly below the chosen
endpoint, and set \(\delta=s^h\). A quadratic branch has
\(l=b+v\) and rate \(c'=b'>0\). A linear branch has
\(l=b_1<b+v\), or is pure linear, and has rate \(c'=b_1'>0\).
The following side constructions are used only on patches where
their strict inequalities hold.

\begin{definition}[Admissible side variable]\label{names:side}
Choose one of the following variables \(Z\).
\begin{enumerate}[label=(\roman*)]
\item \emph{Coarse quadratic side.} Suppose
\(l(a,v)<\min_{[a,D_{\rm end}]}b_1\), and, in the hybrid case,
\(w(a,v)<0\). Fix \(S=\mathcal C(A,V)\) on the same unclipped
branch, with \(A<a\), \(V>v\), and
\[
 V-v\gg(1+c'+|q|+J)(a-A)>0.
\]
Retain all strict domain inequalities and put
\(Z=(S,t_a,R_a,k_a)\). When label recovery is needed for
\(c=1\), also require \(v,V<u(a)\) with slack.
\item \emph{Exact quadratic side, classical.} Let
\(S=(y_0,\bar k,[P^\diamond]_{x_*})\), omitting \(\bar k\)
when absent, with \(x_*>l(a,v)\). The raw coefficient bin must
have zero conditional exponent, or arbitrarily small-cost covering
lists, given \(y_0,t_a\). One may choose
\(x_*<p_*(a)\), and also \(x_*<b_1(D)\) in a virtual
comparison. Put \(Z=(S,t_a,R_a,k_a)\).
\item \emph{Coarse pure linear side.} Take
\(S=\mathcal C(A,V)\), with
\(V-v\gg(1+c')(a-A)>0\), and put \(Z=(S,t_a,L)\).
Maintain \(w<0\) in the hybrid case.
\item \emph{Exact linear side, classical.} Put
\(Z=(y_0,t_a,\bar k,R_a,k_a,L)\), omitting unused entries.
In a clipped branch require \(v>b_1-b\), and remain short of a
virtual endpoint. The additional bins have zero exponent given
\(y_0,t_a\), by ownership and \eqref{eq:A}.
\end{enumerate}
The anchors and static coefficient depths are fixed on a patch.
Countably many such patches cover all the stated admissible roots.
\end{definition}

\begin{lemma}[Common phase]\label{names:common-phase}
Given an admissible side \(Z\), choose a representative point of its
support and let \(C_*\) be its root canonical name. Subtract the
displayed root shear, center and boost in that root cell, and dilate
duration, horizontal width, and normal width by
\(s^a,s^v,s^w\), respectively. In this common frame there is a
bounded phase
\[
 F=(X,B,Y,N,D_0'),\qquad\text{with also }P\text{ on a quadratic branch},
\]
measurable in \(Z\), which differs from the corresponding actual
phase of every conditional observation by \(O(s^\eps)\) for
some positive patch slack. In particular the error is smaller than
every fixed \(\delta\)-power when \(h\) is sufficiently small.
The actual root name and \(C_*\) determine each other to bounded
lists. The sampling scheme and the representative can be chosen
independently of \(v\) on a smaller patch.
\end{lemma}

\begin{proof}
For an exact side, the phase coordinates are fixed by the particle.
In the quadratic case \(x_*>l\) also resolves the residual
derivative to finer precision than the root unit.

For a coarse side first subtract the \(S\)-shear. Its base
uncertainties are \(s^V,s^{A+V},s^{A+2V}\), and its normal
uncertainties are \(s^{w(A,V)},s^{A+w(A,V)}\); the same bounds
hold after flow to \(t_a\). At a representative base point the
root and side derivatives differ by
\(O(s^{\min\{l(a,v),l(A,V)\}})\). Their curvature difference
has size \(O(s^{b(A)})\) in earlier-duration units. Changing to
the fixed representative frame therefore adds velocity error at
most
\[
 O\big(s^{\min\{l(a,v),l(A,V)\}+V}+s^{b(A)+2V}\big),
\]
with an additional factor \(s^A\) for position; omit the second
term in the pure linear case. The inequalities in
Definition~\ref{names:side} put all these errors strictly below
the root widths. The quadratic derivative itself varies by at
most \(O(s^{l(A,V)})\). The change of base tilt fits because
\(A+v+V>a+2v\).

An actual cell may differ from the representative cell at a grid
boundary. Guess its boundedly many neighboring linear-cut and
angular-center cells first. A linear-cut difference multiplies
only unresolved base coordinates at their root widths when
converting residual names; the resulting errors are root normal
widths. The tilt conversion has the same property. This gives
the two-directional lists without comparing normal coordinates
in different frames directly. All fine phase comparisons above
are in the single representative frame.

A measurable representative can be obtained by realizing the
conditional kernel with a uniform seed and fixing a seed that
is valid almost everywhere, by Fubini. The same choice can be
made with the patch parameter under Lebesgue measure, and kept
fixed as \(v\) varies. Only its real coordinates and cell data
are used. Conditional observations still use the original
kernel; no extra shared mark is conditioned upon.
\end{proof}

For the chosen endpoint \(D_{\rm end}\), put
\(T_{\rm end}=T_{D_{\rm end}}\), the start time of its endpoint
interval. On a quadratic branch put
\(\beta_{\rm obs}=\beta_{D_{\rm end}}(T_{D_{\rm end}})\), in
the coordinates of the chosen curvature display. Define the marks,
before common-root dilation of the arguments on the right, by
\begin{equation}
 \theta=\frac{T_{\rm end}-t_a}{s^a},\qquad
 r=\begin{cases}
 \displaystyle\frac{\beta_{\rm obs}-R_a(T_{\rm end})}{s^{b(a)}},
       &\text{quadratic},\\[6pt]
 \displaystyle\frac{P-2R_aX+2k_aB_{t_a}-L}{s^l},
       &\text{linear}.
 \end{cases}
 \tag{xi}\label{eq:xi}
\end{equation}
For the equality case \(c>1\) also include
\(\mathfrak k=(K_1-k_a)s^{a-b(a)}\).
The endpoint is \(1\), or the fixed virtual \(D\). Marks are
bounded. Their joint law with \(Z\) is independent of \(v\) on
the chosen patch; in particular the linear branch has
\(l=b_1(a)\), so its cut \(L\) is independent of \(v\).
Subscripts on marks below denote ordinary bounded-chart bins at
\(\delta\)-depths. Stored marks always have a specified finite
cutoff depth, chosen before the inner scale limit.

\begin{proposition}[Conditional angular dimension]
\label{names:mark-dimension}
Conditional on \(Z\), the time mark has lower surprise
\(dp'\) at every tested positive prefix \(p'\), to arbitrary
strict tolerance. At a typical time root, the joint mark increment
on \(p'\to p'+t\), where \(p'>0\) and \(t\ll p'\), has
conditional surprise at least \(dt\), given its parent and \(Z\).
More precisely:
\begin{enumerate}[label=(\roman*)]
\item If \(c'\geq1\), the time increment itself has this lower
bound given the joint parent. In a linear branch with \(c'>1\),
one may also condition on \(r_{p'+t}\) when
\(p'+t<c'p'\).
\item If \(0<c'<1\), the mixed trace gives \(\nu\in[0,1]\)
such that the \(r\)-increment, given
\(Z,\theta_{p'+t},r_{p'}\), has exponent \(\nu t\), while
the time increment given \(Z,\theta_{p'},r_{p'+t}\) has exponent
at least \((d-c'\nu)_+t\). Their joint increment has exponent
at least \(dt\).
\end{enumerate}
In the classical exact linear case with \(Q\leq1\) and
\(0<b_1'<1\), \(\nu\) is the raw final derivative trace with
side \(y_0\) and graph \(b_1\); this also applies strictly
before a virtual endpoint with \(b_1'>0\).
\end{proposition}

\begin{proof}
At the gate \(i=a+p'h\), every predecessor on a tested completing
path predicts the static side and the other root data to small
lists. For a coarse side this uses exact base flow, the strict
normal slack for hybrid backflow, and \eqref{eq:A}. For an
exact predicted coefficient, use its arbitrarily small-cost list
and discard its negligible final failure event. The row predicts
\(\theta_{p'}\), and it predicts \(r\) to depth \(p'\) if
\(c'\geq1\), or depth \(c'p'\) otherwise. In a linear branch
the latter accuracy always holds. When \(c>1\), the predicted
curvature trajectory is evaluated at the endpoint; its bounded
normalized slope makes replacing that time by the gate time
cost only \(\delta^{p'-o(1)}\).

Here is the causal counting step. Use Lemma~\ref{cyl:borel-paths}:
take a Borel inner version of the prediction-success event, mark
rows at a later gate \(j'\) that have a continuation to it, and
retain a Borel inner version of that analytic mark under the full
analyzed \(j'\)-row energy law, inside the current end-state support.
For each of their times, restrict
the \(i\)-inputs to
a Borel outer version of the listed predecessor set. It includes all
predecessors under the unrestricted truncated \(i\)-to-\(j'\)
geometric correspondence, not just predecessors on successful
completions, and has the same full analyzed input energy. A path from a
complementary input to a marked row would
concatenate with its completing path, and
so would contradict the definition of that predecessor list.
Correspondence locality therefore makes complementary transfer
negligible. Ordinary contraction, summed over the listed next
time cells, costs their multiplicity times the raw mass at
\(i\). Applying \eqref{eq:Sp} onward to the readout prohibits
lists of exponent strictly below \(d(j'-i)\), including the
small predictor multiplicities. Testing popular cells proves
the stated time lower surprises. The same proof starts at
\(p'=0\).

For \(c'<1\), use Lemma~\ref{names:mixed-trace} with
\(V_*=\beta^\diamond,B=b\) for constant quadratic cuts, or
with the raw slope and \(B=b_1\) in a linear branch. The side
is the specified static side. Its additional root data have
zero exponent beyond \(S,t_a,[V_*]_{B(a)}\). At the later
time, \(p'>c'(p'+t)\) for sufficiently small \(t/p'\), so
the \(r\)-increment lies above the graph and has rate \(\nu\).
Adding the coefficient bins between \(B(i)\) and the tested
depth to the causal time bound loses at most
\(c'(\nu-\nu_-)t\): this is precisely the mixed coefficient
information exposed by the new time bits. The \(r\)-increment
before those bits has rate at least \(\nu\). Chain rule gives
joint rate at least
\[
 d-c'(\nu-\nu_-)+\nu\geq d,
\]
and also the separately stated time bound. Known shifts of bin
origins and virtual raw coefficient comparisons cost small
lists only.
\end{proof}

\subsection{Separation forced by a positive minimal cost}

\begin{proposition}[Angular separation]\label{names:separation}
In an equality quadratic family with \(c>0\), the conditional
\((\theta,r)\) law for an admissible side avoids arbitrary
affine lines with a positive power. At \(c=1\) this assertion
requires the coarse side with the additional left label slack
in Definition~\ref{names:side}; it is not asserted for a general
exact side. More precisely, for a line chosen measurably per
\(Z\), its \(\delta^\lambda\)-neighborhood has total
probability at most \(\delta^{\eps\lambda}\) in the exponent
limits, for some \(\eps>0\), each tested \(\lambda>0\), and
sufficiently small \(h\). The constants are uniform on bounded
lists of positive \(\lambda\).

In a pure linear family, the analogous statement holds for
predicting a value of \(r\). In a clipped classical linear
branch with \(p_*(a)=b_1(a)\), two conditionally independent
marks satisfy
\(\abs{r_i-r_j}\geq\delta^\lambda\) with probability tending
to one for every fixed \(\lambda>0\); no uniform positive
avoidance exponent is asserted in this case.
\end{proposition}

\begin{proof}
For the clipped assertion, root additions to \(y_0,t_a\) have
zero cost and the offset in \eqref{eq:xi} is then known. Every
strictly finer raw slope bin has positive deterministic score
by the definition of \(p_*\). After trimming exceptional heavy
bins, conditional independent sampling gives the separation.
In a virtual use keep the query below \(b_1(D)\).

We prove the power statements by contradiction. Failure gives
a predicted tube of width \(s^{(1-o(1))z}\), where
\(z=\lambda h>0\), carrying mass \(s^{o(z)}\); all previous
errors are taken to zero in \(z\)-units first. A nearly vertical
line confines the time coordinate to a power-small interval and
is excluded by Proposition~\ref{names:mark-dimension}. Hence,
after an arbitrarily small loss, the line is the graph of an
affine function with slope \(s^{-o(z)}\) in mark coordinates.
In physical units it predicts an affine trajectory \(\beta_*\)
with
\[
 |\beta_1(T_1)-\beta_*(T_1)|
       \leq s^{b(a)+(1-o(1))z}.
\]
Exact extra-prediction sides are restricted to their successful
lists, taking their costs and failure tolerances sufficiently
small in \(z\)-units.

We shall extract a segment \([i,j]\). By
Lemma~\ref{cyl:borel-paths}, take a Borel inner version of the
retained final event, mark \(j\)-rows having a completion to it,
and then take a Borel inner version of this analytic mark under the
full analyzed \(j\)-row energy law, inside the current end-state
support. Every retained marked row still has a
completion. Their mass is at least
\(s^{o(z)}E_j\) by \eqref{eq:Sp}; the count bound gives a
cube-sum lower bound \(s^{o(z)}E_j\sqrt{\kappa_j}\).
Group starting rows by the binned polynomial subtraction and
their starting interval. Disjointize the finite possible-bin
relations and retain their Borel inner versions, so each marked end
row has a bin witnessed by a completion. Use the Borel outer
predecessor set for that bin, formed from the unrestricted truncated
\(i\)-to-\(j\) geometric correspondence. It includes all its
correspondence predecessors without changing full analyzed input energy;
its Borel inner partner carries the same full analyzed state inside
the original starting-bin partition. Each end label has
\(s^{-o(z)}\) chosen bins, as checked in the cases below.
Path concatenation and locality
then justify restricting synthesis to those starting sets.
Input energy summed over groups is at most
\(s^{-o(z)}E_i\), each group's regularity is at most
\(s^{-o(z)}\kappa_i\), and the end counts gain only the same
multiplicity. H\"older preserves the cube-sum lower bound
over groups. After rescaling, these grouped experiments realize
the preceding extremal parameters in the closure. Thus a fixed
improvement in the minimal cost is impossible.

Suppose first \(0<c<1\). Set
\(i=a+z/2\), \(j=a+3z/4\), and choose
\(0<\omega<\min\{(1-c)/2,c/4\}\). Subtract the prediction
at the \(i\)-time start, rounded at depth \(b(i)+\omega z\).
Starting rows know it from the side. Moving the prediction to
the endpoint time costs at most
\(s^{b(a)+z/2-o(z)}\), and comparing the end label with the
final coefficient costs \(s^{b(j)-o(z)}\). Both are finer
than the chosen bin, so the end label knows few bins.
The residual end coefficient is at most
\(s^{b(i)+\omega z-o(z)}\), while the duration-scaled own
slope already has size \(s^{b(j)-o(z)}\). In segment units
the curvature cost is at most \(c-4\omega+o(1)<c\).

For \(c>1\), set \(i=a\), \(j=a+\gamma z\), with
\(1/c<\gamma<1\), and choose \(\gamma<\tau<1\).
On the part where
\(s^j|K_1-\beta_*'|>s^{b(a)+\tau z}\), the end trajectory
has the same slope discrepancy to within a factor two, since
\(c\gamma>1>\tau\). Its meeting the predicted trajectory
within the tube confines the terminal time, per end row and
predicted line, to length
\(s^{j+(1-\tau-o(1))z}\). The causal time count excludes
this part by a power. On the remainder, extrapolation to
\(t_a\) gives coefficient and slope errors bounded by
\[
 s^{b(a)+(\tau-\gamma)z-o(z)},\qquad
 s^{b(a)-a+(\tau-\gamma)z-o(z)}.
\]
Subtract the predicted pair at those base depths improved by
\(\omega z\), with \(0<\omega<\tau-\gamma\).
It is jointly list-recoverable, and both residual curvature
terms at the end have cost strictly below \(c\).

For \(c=1\), use \(i=a\), \(j=a+2z\). The additional
left slack makes \(\lambda_j\) predict the coarse side: its
base width has \(u(j)>V\), its derivative comparison fits
\(l(A,V)\), and its normal comparison fits by the same
Taylor and hybrid-drift bounds as in
Proposition~\ref{names:canonical-entropy}. It also predicts
the root curvature bins. Subtract the predicted pair at
depths \(b(a)+\omega z,b(a)-a+\omega z\), where
\(0<\omega<1/2\). The residual end coefficient is bounded
by \(s^{b(a)+\omega z-o(z)}\); the end-duration times the
residual slope is bounded by \(s^{b(j)-o(z)}\). Hence the
end curvature cost is strictly below one. This case does not
require the slope of the prediction to be inferred from the
end curvature alone: it is predicted by the recovered side.

Finally, in the pure linear case choose \(i=a\),
\(j=a+\gamma z\) with \(c_1\gamma>1\). A value prediction
for \(r\) predicts the terminal linear coefficient at depth
\(b_1(a)+z-o(z)\). Subtract its bin at
\(b_1(a)+\omega z\), \(0<\omega<1\). The end label knows
the bin by its finer coefficient comparison. The curvature
cost remains zero and the linear cost becomes at most
\(c_1-\omega/\gamma+o(1)\), a contradiction.

For clarity, the recovery margins behind the group count are as
follows. When \(0<c<1\), the prediction-transport and end-label
errors have depths \(b(a)+z/2\) and
\(b(j)=b(a)+3cz/4\), both finer than the subtraction depth
\(b(i)+\omega z=b(a)+(c/2+\omega)z\) by the two restrictions on
\(\omega\). When \(c>1\), both errors in the displayed value--slope
pair exceed their subtraction depths by
\((\tau-\gamma-\omega)z>0\). At \(c=1\), the recovered coarse
side supplies the slope bin to a small list, and the end coefficient
has depth \(b(j)=b(a)+2z>b(a)+\omega z\). In the pure linear case
the terminal coefficient has depth
\(b_1(j)=b_1(a)+c_1\gamma z>b_1(a)+\omega z\).
Thus the pre-existing side and rounding lists leave only
\(s^{-o(z)}\) bins per fixed end label in each case.
\end{proof}

\subsection{Replica budgets and conditional fibers}

Draw finitely many observations independently given \(Z\), and
store their marks \(\Xi_i=(\theta_i,r_i)\), including
\(\mathfrak k_i\) when required, to fixed finite depths.
Let \(F=F(Z)\) and \(C_*=C_*(Z)\) be the representative phase
and name from Lemma~\ref{names:common-phase}. The cutoff for
each stored mark is chosen sufficiently large for the finitely
many ensuing queries; it is never replaced by an exact mark
inside an information estimate without a new justification.

For the normalized phase and angle, \(F_{\leq p'}\) and
\(X_{\leq p'}\) denote their coordinatewise dyadic names at
\(\delta\)-depth \(p'\); these names also determine their coarser
dyadic bins.

\begin{proposition}[Replica information budgets]
\label{names:replica-budgets}
At almost every interior admissible root there are arbitrarily
small tangent gaps with sufficiently accurate preceding
experiments such that, for any fixed finite collection of
queries,
\begin{equation}
 I_\delta(F_{\leq p'};\boldsymbol\Xi\mid C_*)=o(1).
 \tag{vb}\label{eq:vb}
\end{equation}
Define a fiber read \(Q_i\), at a sufficiently large finite
depth relative to still finer stored marks, by binning
\[
 B+\theta_iX,\quad Y+\theta_iX^2,\quad D_0'+\theta_iN,
 \quad
 \begin{cases}
 P-2r_iX,\ N-PX+r_iX^2,&\text{quadratic},\\
 N-r_iX,&\text{linear}.
 \end{cases}
\]
Then
\begin{equation}
 \begin{split}
 I_\delta(X_{\leq p'};Q_i,\boldsymbol\Xi\mid C_*)&=o(1),\\
 i_\delta(X_{\leq p'}\mid C_*)&=m(a,v)p'+o(1)
                         \quad\text{in probability}.
 \end{split}
 \tag{fib}\label{eq:fib}
\end{equation}
At a fixed boundary root with one observation, only the
consequence of \eqref{eq:ng} for adding that observation and
its fiber read is asserted; neither \eqref{eq:vb} nor an
interior angular trace is asserted there.
\end{proposition}

\begin{proof}
The marks and their law before evaluation of the representative
phase are independent of \(v\) on the patch. Their total
entropy is \(O(h\log(1/s))\). For sufficiently large fixed
\(K\), the pair \(C_*(v),C_*(v+Kh)\) determines
\(F_{\leq p'}\) up to bounded lists. In a quadratic template
the new derivative bin reads \(P\) to depth \(K\). The new
residual phase can be converted to the old shear frame using
the known shear difference and finer base bins; every old
coordinate is then known to depth at least \(K\). In a
linear template the linear cut is unchanged. Consequently
telescoping the decreasing function
\(H_s(\boldsymbol\Xi\mid C_*(v))\) over a translation
of size \(Kh\) bounds the integral of
\(I_s(F_{\leq p'};\boldsymbol\Xi\mid C_*)\) by
\(O(h^2)+o(1)\). Here \(I_s=hI_\delta\), since
\(\log(1/\delta)=h\log(1/s)\).

Integrate also in \(a\) on compact patches when needed.
For fixed dyadic \(h\), take lower limits of the sum of the
finitely many nonnegative errors through the earlier
approximations. Fatou preserves the integral bound. After
division by \(h\), the dyadic integrals are summable, so
almost every root admits accurate experiments with all the
required normalized errors tending to zero. Countably many
finite cutoffs and tuple sizes are handled by a diagonal.
This establishes existence of the stated tangent laws; it
does not assert stationarity for arbitrary sampling schemes
that change with \(v\).

Consider next observation \(i\) and its actual child name
\(\mathcal C(a+Kh,v)\), with \(K\) larger than all
required query depths, also after multiplication by \(c'\)
and, for nonconstant quadratic cuts, by \(c'-1\).
It predicts \(C_*\) and reads \(\Xi_i,Q_i\) to bounded
lists. Here are the details for the last assertion.
The child time \(\theta'\) approximates \(\theta_i\)
arbitrarily finely. Write
\(W'=B+\theta'X\), \(V'=Y+\theta'X^2\).
The difference between child and root shears, in the common
frame, is
\[
\begin{split}
 c_{\rm obs}+\mathcal LX+\mathcal RX^2
                 +\mathfrak b(V'-2XW'),\\
 d_{\rm obs}+\mathcal LW'+\mathcal RV'
                 -\mathfrak b(W')^2.
\end{split}
\]
All coefficients and offsets are determined by the two names.
In the quadratic case \(\mathcal R,\mathfrak b\) approximate
\(r_i,\mathfrak k_i\), and
\(\mathcal L+2\mathcal RX-2\mathfrak bW'\) approximates
\(P\). Thus \(\mathcal L-2\mathfrak bW'\) reads
\(P-2r_iX\); adding back the known
\(\mathfrak bV'\) to the residual velocity reads
\(N-PX+r_iX^2\). The base reads supply \(W',V'\), and
the residual position supplies \(D_0'+\theta' N\).
For a linear branch the same derivative expression reads
\(r_i\), while \(\mathcal R,\mathfrak b\) are negligible
because \(b+v>l\), or identically zero. It therefore reads
\(N-r_iX\). Changing the curvature basis in a clipped
child costs \(s^{b+v}\), smaller than every required
\(s^l\delta^{M'}\) by the fixed branch slack.
Nonconstant quadratic cuts read \(r_i\) with error
\(O(\delta^K+\delta^{c'K})\) and the slope to depth
\((c'-1)K\). Large constant offsets are converted at the
exact child-bin time before comparison. Finally the common
phase error is negligible by Lemma~\ref{names:common-phase}.

Equation~\eqref{eq:ng} bounds the information that the child
name gives about the fine angle, conditionally on the root.
Thus adding \((Q_i,\Xi_i)\) costs \(o(1)\).
To add the other marks, first include a sufficiently fine
phase bin. Given it and \(\Xi_i\), the fiber read has
only bounded ambiguity. The remaining conditional phase
information is bounded by \eqref{eq:vb}, using chain rule.
This proves the first line of \eqref{eq:fib}. The second
line is the deterministic angular profile and the root
comparison from Lemma~\ref{names:common-phase}.
\end{proof}

\begin{corollary}[Stationary conditional walks]\label{names:walks}
One may resample the common side and its phase successively
from their joint law conditional on
\(C_*,\boldsymbol\Xi,Q_i\), optionally retaining an angular
prefix bin. Each leg preserves the joint marked marginal.
Fine angular increments of the new endpoint have negligible
normalized information from the preceding history beyond
that already present in \(C_*\) and the retained prefix.
For a leg with \(A=\Delta X\), its increments satisfy, up
to arbitrarily fine specified errors,
\begin{equation}
\begin{aligned}
 \Delta B&=-\theta_i A,&
 \Delta Y&=-\theta_i(2XA+A^2),&
 \Delta D_0'&=-\theta_i\Delta N,\\
 \Delta P&=2r_iA,&
 \Delta N&=PA+r_iA^2&&\text{(quadratic)},\\
 &&\Delta N&=r_iA&&\text{(linear)}.
\end{aligned}
\tag{leg}\label{eq:leg}
\end{equation}
For a finite walk without a retained prefix, the new angular
draws have joint negligible information relative to independent
draws given \(C_*\), including relative to the starting
phase and all marks.
\end{corollary}

\begin{proof}
Conditional resampling preserves the conditioning data and
their joint marginal. Given those data it is independent of
the rest of the history. Apply \eqref{eq:fib} to this
unchanged marginal; chain rule shows that retaining an
angular prefix only decreases the information budget for
the remaining angle bits. Summing the finite number of
budgets gives the joint assertion. The first three
identities in \eqref{eq:leg} follow from equality of the
first three fiber reads. The derivative read gives
\(\Delta P=2r_iA\). Substitution in the remaining quadratic
read gives \(\Delta N=PA+r_iA^2\); the linear read gives
the last identity. Increase the finite cutoff to make all
rounding errors smaller than the queried scales.
\end{proof}

These walk assertions are information estimates, not
total-variation assertions. In particular they do not permit
retaining a set of merely subpower probability and reusing
the old budget without a separate argument.

\subsection{Random scores and diagonal traces}
\label{names:score-calculus}

The walk estimates provide small information budgets, while the projection
tests to come use scores of individual outcomes at several depths. We next
justify passing from finite-name scores to local rates and specify which
changes of probability law preserve those budgets.

Throughout this subsection, $\rho\downarrow0$ and scores are divided by
$L_\rho=\log(1/\rho)$.  Side data $D$ may take values in an arbitrary
standard Borel space.  All conditional laws are regular conditional laws.
For a finite name $A$, write
\[
 i_\rho(A\mid D)
   =-L_\rho^{-1}\log\Prob(A=A(\omega)\mid D).
\]
Every assertion about several names is made first for a fixed finite list
of queries.  Countable limits below are limits of these joint score laws;
they do not assert a coupling of experiments at different scales.

\begin{lemma}[Pathwise score limits]
\label{names:score-limits}
Suppose a finite collection of name arrays $A_q$, with depths in a bounded
interval, has the following properties.  Coarser names are determined
from finer names up to lists of size $\rho^{-o(1)}$, and, conditionally on
$A_q,D$, the name $A_{q'}$, $q'\ge q$, has at most
$\rho^{-K(q'-q)-o(1)}$ possible values.  The same bounds hold for the
joint arrays under consideration, and their total range sizes are at
most a fixed power of $\rho^{-1}$.  Then, after a subsequence, the scores
at countably many fixed queries have a joint distributional limit with
these properties:
\begin{enumerate}[label=\textup{(\roman*)}]
\item Cumulative scores are nonnegative, nondecreasing, and Lipschitz in
each named-variable refinement argument, with conditioning held fixed
and with the corresponding branching constant.
They extend from rational queries to real queries.
\item The chain rule holds, and adding conditioning can only decrease a
conditional score.  In particular, for $r\le r'\le a\le b$, whenever
$F_r$ is a nested conditioning array,
\begin{align}
 &i(A_b\mid F_{r'},D)-i(A_a\mid F_{r'},D)\notag\\
 &\hspace{15mm}\le
 i(A_b\mid F_r,D)-i(A_a\mid F_r,D).
 \label{names:score-increment-monotonicity}
\end{align}
Here and below $i$ without a scale denotes a limiting random score.
\item Mutual determination with exponent-zero lists identifies the
corresponding limiting scores.  All these statements hold simultaneously
almost surely for the countable family of queries.
\end{enumerate}
\end{lemma}

\begin{proof}
If a conditional distribution is supported on a list of size $N$, then
the total probability of values of conditional probability less than
$\rho^\epsilon/N$ is at most $\rho^\epsilon$.  Apply this observation to
each refinement and each list comparison.  It bounds the refinement
score between zero and $K(q'-q)+\epsilon+o(1)$, except on a set of
probability tending to zero.  Bounded total range gives the same
control on whole scores; values exceeding the range exponent by
$\epsilon$ have probability at most $\rho^\epsilon$, also conditionally
on $D$.  Thus every finite vector of scores is tight.  A diagonal
subsequence gives a consistent joint limit on the countable query set.

For exactly nested names the log-probability chain rule is an identity.
For list nesting, insert the coarser name into the finer name.  Its
additional conditional score tends to zero by the preceding list bound,
so the same identities survive in the limit.  Equation~\eqref{eq:LP}
gives the conditioning inequality at each fixed query.  Apply it to the
increment name with the coarse name already conditioned upon to obtain
Equation~\eqref{names:score-increment-monotonicity}.  A countable union
of null sets gives all inequalities simultaneously.  The loss from
refining a conditioning name $F_r$ to $F_{r'}$ is nonnegative and at
most $i(F_{r'}\mid F_r,D)$, by the chain rule and nonnegativity of
conditional scores.  It is therefore at most $K(r'-r)$ with the
appropriate branching constant.  Thus the required mixed score
functions are Lipschitz also in their conditioning depths.  The Lipschitz
estimates extend the limiting scores continuously from rational depths.
Comparisons with neighboring rational depths also give the claimed
limit at any additional fixed real query.  The same reasoning applies
when several arrays are refined together.
\end{proof}

The following elementary trace statement is useful because differentiating
a two-variable Lipschitz function on its diagonal is not justified by
the usual almost-everywhere differentiability theorem alone.

\begin{lemma}[Monotone diagonal trace]
\label{names:score-trace}
Let $H(q,r)$ be continuous for $0<r\le q<T$, with $H(r,r)=0$.
Assume that $H(\cdot,r)$ is nondecreasing and $K$-Lipschitz, uniformly in
$r$, and that
\[
 H(b,r')-H(a,r')\le H(b,r)-H(a,r)
 \qquad(r\le r'\le a\le b).
\]
There is a measurable function $g:(0,T)\to[0,K]$ such that, for almost
every $p$, uniformly over $p\le u\le v\le p+t$,
\begin{equation}
 H(v,u)=g(p)(v-u)+o_p(t)
 \qquad(t\downarrow0).
 \label{names:score-triangular-trace}
\end{equation}
\end{lemma}

\begin{proof}
For rational $r$, choose the derivative
$f_r(q)=\partial_qH(q,r)$ outside one common null set.  Increment
monotonicity gives $0\le f_{r'}(q)\le f_r(q)\le K$ for rational
$r<r'<q$.  Put
\[
 g(q)=\inf_{\substack{r\in\mathbb Q\\0<r<q}}f_r(q).
\]
For any fixed real $u$, differentiate the increment inequalities
comparing $H(\cdot,u)$ with all rational conditioning depths.  For
almost every $q>u$ and each rational $r<u$, they imply
\[
 g(q)\le\partial_qH(q,u)\le f_r(q).
\]
For the first inequality, compare with a rational depth between $u$
and $q$.  Absolute continuity therefore gives, for every $u<v$ and
rational $r<u$,
\begin{equation}
 \int_u^v g(q)\,dq\le H(v,u)\le\int_u^v f_r(q)\,dq.
 \label{names:score-trace-sandwich}
\end{equation}
Choose $p$ to be a Lebesgue point of $g$ and of every $f_r$ defined in a
neighborhood of $p$, and outside the common derivative null set.  These
conditions exclude only a null set.  For fixed rational $r<p$, the lower
and upper errors in Equation~\eqref{names:score-triangular-trace} are
bounded, uniformly over the indicated triangle, by
\[
 \int_p^{p+t}|g(q)-g(p)|\,dq
 \quad\hbox{and}\quad
 \int_p^{p+t}|g(q)-g(p)|\,dq
       +\int_p^{p+t}(f_r(q)-g(q))\,dq,
\]
respectively.  The first is $o(t)$.  The limsup of the second divided by
$t$ is at most $f_r(p)-g(p)$.  Taking the infimum over rational $r<p$
makes it zero and proves the assertion.
\end{proof}

\begin{corollary}[Prefixes with future conditioning]
\label{names:score-future-prefix}
Let $F_q,U_q,V_q$ satisfy Lemma~\ref{names:score-limits}, and suppose
$F_q,D$ determines $U_q,V_q$ to exponent-zero lists.  For almost every
fixed $p$, almost surely in the limiting score law, there are rates
$g_U(p),g_V(p),g_{UV}(p)$ such that
\[
 i(U_{p+t}\mid F_p,D)=g_U(p)t+o(t),\qquad
 i((U,V)_{p+t}\mid F_p,D)=g_{UV}(p)t+o(t).
\]
Moreover, for every fixed $0\le\lambda\le1$,
\begin{equation}
 i(U_{p+\lambda t}\mid F_p,V_{p+t},D)
  =\lambda\bigl(g_{UV}(p)-g_V(p)\bigr)t+o(t).
 \label{names:score-future-prefix-formula}
\end{equation}
The error can be taken uniform in $\lambda$.  Tuples may replace $V$.
\end{corollary}

\begin{proof}
For $A=U,V,(U,V)$, apply Lemma~\ref{names:score-trace} pathwise to
$H_A(q,r)=i(A_q\mid F_r,D)$.  Its diagonal vanishes by the list
hypothesis, and its increment inequality is
Equation~\eqref{names:score-increment-monotonicity}.  Put
$u=p+\lambda t$ and $v=p+t$.  The chain rule first gives
\[
 i(U_u\mid F_p,V_u,D)
   =H_{UV}(u,p)-H_V(u,p)
   =\lambda(g_{UV}(p)-g_V(p))t+o(t).
\]
Write $i(A;B\mid C)=i(B\mid C)-i(B\mid A,C)$ for limiting conditional
information.  These quantities are nonnegative by
Equation~\eqref{eq:LP}.  Since $U_u$ is determined by $F_u,D$, the
additional loss from the future part of $V$ satisfies
\begin{align*}
 0&\le i(U_u;V_v\mid F_p,V_u,D)\\
  &\le i(F_u;V_v\mid F_p,V_u,D)\\
  &=H_V(v,p)-H_V(u,p)-H_V(v,u)=o(t).
\end{align*}
The second inequality follows by expanding the information involving
$(U_u,F_u)$ in both orders and using nonnegativity.  The last equality
uses nesting and determination of $F_p,V_u$ from $F_u,D$; its $o(t)$
bound follows from the uniform triangular trace.  This proves
Equation~\eqref{names:score-future-prefix-formula}.  Finally, Fubini's
theorem converts the pathwise almost-everywhere statement into an
almost-sure statement at almost every fixed $p$.
\end{proof}

\begin{lemma}[Changes of probability law]
\label{names:score-law-change}
Let $P_\rho,Q_\rho$ be probability laws with
$\operatorname{KL}(P_\rho\Vert Q_\rho)=o(L_\rho)$.  At samples drawn
from $P_\rho$, their scores for any fixed finite collection of names,
including conditional scores with arbitrary common side data, differ
by $o(1)$ in probability.  If $P_\rho(E)\ge c>0$, then
$Q_\rho(E)\ge\rho^{o(1)}$.

Restricting a law to an event of probability at least $c>0$ preserves
its finite-name scores up to $o(1)$ in probability.  It also preserves
any conditional mutual-information bound $o(L_\rho)$ between two
variables given common side data.  The latter assertion is not made
for events whose probabilities merely have exponent zero.
\end{lemma}

\begin{proof}
Let $R=dP_\rho/dQ_\rho$.  The elementary estimate
$P_\rho(\log R<-u)\le e^{-u}$ implies
$\E_{P_\rho}(\log R)_-\le1$, and hence
\[
 \E_{P_\rho}(\log R)_+
   \le\operatorname{KL}(P_\rho\Vert Q_\rho)+1=o(L_\rho).
\]
Markov's inequality proves $\log R=o(L_\rho)$ in probability.
Relative entropy decreases under measurable maps, so this also applies
to the likelihood ratios for each joint name and its conditioning
field.  Subtracting those two log ratios proves the conditional-score
assertion.  Applying the log-sum inequality to $E,E^c$ gives
\[
 \operatorname{KL}(P_\rho\Vert Q_\rho)
 \ge P_\rho(E)\log\frac1{Q_\rho(E)}-\log2,
\]
which proves the event assertion.

For restriction to $E$, use
$\operatorname{KL}(P_\rho(\cdot\mid E)\Vert P_\rho)=\log(1/P_\rho(E))$.
To verify the information assertion, let $X,W,D$ denote the variables
in question and set $Q=P_D P_{X\mid D}P_{W\mid D}$.  With $c_\rho=P(E)$,
the mixture identity for relative entropy gives
\[
 c_\rho\operatorname{KL}(P_{XWD\mid E}\Vert Q)
 \le I_P(X;W\mid D)+h(c_\rho),
\]
where $h(c)=-c\log c-(1-c)\log(1-c)\le\log2$.
This identity remains valid if $E$ also depends on auxiliary variables:
apply the mixture identity to the marginal laws of $(X,W,D)$.
The relative entropy on the left is at least
$I_{P(\cdot\mid E)}(X;W\mid D)$, by decomposition against the product
of the restricted conditional marginals.  Division by $c_\rho\ge c$
proves the result.
\end{proof}

These results are used without selecting rare limiting profiles.  A
strict violation is tested at a fixed admissible $p$, on an event of
positive probability in its random score law, using finitely many
queries and slightly relaxed tolerances.  The joint convergence of
those scores transfers positive probability to sufficiently advanced
experiments.  Their information errors are then taken sufficiently
small for this fixed event and layer, before the layer shrinks.

\begin{remark}[Freezing coefficients]
\label{names:score-coefficients}
Suppose a local projection has smooth coefficients depending on the
phase point, such as the contact differential $dY-2X\,dB$.  Fix their
values to accuracy $O(\rho^{p_0})$ using a name already in the
conditioning, where $0<p_0<p$.  Inside an $F_p$ cell, freezing changes
an increment by $O(\rho^{p+p_0})$; Taylor expansion of the original
coordinate contributes $O(\rho^{2p})$.  Both are smaller than
$\rho^{p+t}$ when $0<t<\min(p,p_0)$.  The two descriptions therefore
give bounded lists for one another at this layer, conditionally on the
parent and the coefficient descriptions.  The values at the parent
representative must be translated exactly: they are known offsets and
are not increment errors.  The same argument applies to a mark
coefficient frozen to accuracy $\rho^\epsilon$ when $t<\epsilon$.
Any cost of adding this coefficient information is a separate
information budget; the geometric comparison does not supply it.

For a Heisenberg recentering, the error left by changing horizontal
cell representatives is bilinear in unresolved horizontal
displacements.  If their scale is $\rho^a$, this error is
$O(\rho^{2a})$.  Thus a vertical name at scale $\rho^b$ is predicted
up to bounded lists whenever $b\le2a$, giving the required nesting and
branching in that wedge.  The score calculus applies in its interior;
it supplies no derivative trace across the boundary $b=2a$.
\end{remark}

\subsection{A projection rule with conditional records}
\label{names:score-projections}

The record in the next proposition may contain more information than
the direction it determines.  In particular, different records with
the same direction may carry different lists of projection values.

\begin{proposition}[Conditional projection rule]
\label{names:projection-rule}
Let $0\le S\le2$, $0<b\le1$, and $\Delta\downarrow0$.  Let $Z$ be a
bounded planar source, $W$ a record determining a direction $g(W)$,
and $D$ common side data.  Assume
\begin{equation}
 I(Z;W\mid D)=o(\log(1/\Delta)).
 \label{names:score-projection-budget}
\end{equation}
Suppose there is an event $G$ of probability bounded below such that
its unnormalized conditional marginals obey, for every needed ball
and interval of radius $\Delta\le r\le1$,
\begin{align}
 \Prob(G,\ Z\in B(z,r)\mid D)&\le\Delta^{-o(1)}r^S,\notag\\
 \Prob(G,\ g(W)\in J(r)\mid D)&\le\Delta^{-o(1)}r^b.
 \label{names:score-projection-masses}
\end{align}
For each $D,W$, let $\mathcal I(D,W)$ be a list of at most
$\Delta^{-x-o(1)}$ intervals of length $O(\Delta)$, and suppose that on
$G$ the projection of $Z$ in direction $g(W)$ belongs to their union,
up to $O(\Delta)$ error.  Then
\begin{equation}
 x\ge\min\{S,(S+b)/2,1\}.
 \label{names:score-projection-conclusion}
\end{equation}
The assertion is a strict-error statement: a fixed strict deficit is
impossible if all preceding exponent losses and the information ratio
are sufficiently small.  Consequently a sufficiently fine finite grid
of prefix radii suffices, and errors of size
$\Delta^{1-o(1)}$ in directions, representatives, or output resolution
are allowed.
The same conclusion holds with $c_\Delta=\Prob(G)$ tending to zero,
provided
\[
 \log(1/c_\Delta)=o(\log(1/\Delta)),\qquad
 I(Z;W\mid D)=o(c_\Delta\log(1/\Delta)).
\]
\end{proposition}

\begin{proof}
The case $S=0$ follows from the nonempty output list on $G$, so assume
$S>0$.  A finite choice of rotated projective charts allows restriction
to a positive-probability part of $G$ where projections have the form
\[
 \pi_u(z)=z_1+uz_2,
\]
with bounded $u,z$.  Multiplying a projection by a bounded nonzero
factor only changes interval constants.  First restrict to $G$.
Lemma~\ref{names:score-law-change} preserves the vanishing normalized
information budget.  If $c_D=\Prob(G\mid D)$ and $\Prob(G)\ge c$,
then under this restricted law the parents with $c_D<c/2$ have total
probability at most $1/2$.  On the other parents, normalization of
Equation~\eqref{names:score-projection-masses} costs at most $2/c$.
The conditional information is nonnegative and has vanishing average
after normalization by $\log(1/\Delta)$.  We can therefore choose a
parent with these mass bounds and with conditional information
$o(\log(1/\Delta))$.

Fix that parent.  Write $P$ for the restricted joint law and
$\mu,\nu$ for its source and record marginals.  Let $R$ be the relation
that the projected source is in the record's output list.  Since
$P(R)=1$, the log-sum inequality gives
\[
 \eta:=(\mu\times\nu)(R)
   \ge\exp\bigl(-\operatorname{KL}(P\Vert\mu\times\nu)\bigr)
   =\Delta^{o(1)}.
\]
Thus product sampling still hits the lists, although a fixed positive
product probability is not asserted.

Fix a small $\epsilon>0$.  For sufficiently advanced experiments,
the marginal bounds have constants at most $\Delta^{-\epsilon}$ and
$\eta\ge\Delta^\epsilon$.  Draw
$M=\lceil\Delta^{-b}\rceil$ independent full records
$W_1,\ldots,W_M$ from $\nu$.  With probability tending to one faster
than any fixed power of $\Delta$, their empirical direction law
$\nu_M$ obeys
\begin{equation}
 \nu_M(J(r))\le C\Delta^{-2\epsilon}r^b
 \qquad(\Delta\le r\le1).
 \label{names:score-empirical-directions}
\end{equation}
Indeed, for each dyadic interval of radius $r$, the count is binomial
with mean at most $M\Delta^{-\epsilon}r^b$; the asserted threshold is
larger by a factor comparable to $\Delta^{-\epsilon}$ and is at least
a constant multiple of $\Delta^{-2\epsilon}$.  The binomial upper-tail
bound is $\Prob(N\ge k)\le(eMq/k)^k$ for a count with $M$ trials and
success probability $q$: sum over the $k$ successful indices and use
$\binom{M}{k}\le(eM/k)^k$.  This bound, followed by a union bound over
$O(\Delta^{-1})$ dyadic intervals down to scale $\Delta$, proves the
claim.  Every other interval is
covered by a bounded number of these intervals at comparable radius.

The expected average relation density of this sample is $\eta$.
Since the probability that
Equation~\eqref{names:score-empirical-directions} fails is $o(\eta)$,
there is a sample satisfying that equation and
\[
 \int\theta(z)\,d\mu(z)\ge\eta/2,
 \qquad
 \theta(z)=M^{-1}\#\{j:(z,W_j)\in R\}.
\]
Let $E=\{z:\theta(z)\ge\eta/4\}$.  The bounds $0\le\theta\le1$
imply $\mu(E)\ge\eta/4$.  The probability
$\mu_E=\mu|_E/\mu(E)$ has planar ball masses at most
$C\Delta^{-2\epsilon}r^S$.  For each $z\in E$, put the uniform
probability on its successful sampled records.  Its direction masses
are at most $C\Delta^{-3\epsilon}r^b$, by
Equation~\eqref{names:score-empirical-directions} and
$\theta(z)\ge\eta/4$.

Dualize: the source point $z$ parametrizes the line
\[
 \ell_z=\{(u,z_1+uz_2):u\in\R\}.
\]
The map from $z$ to the line parameters is bi-Lipschitz on the bounded
chart, so this family supports an $S$-Frostman probability with the
preceding small power loss.  On each $\ell_z$, the successful records
give a $b$-Frostman probability: the direction parameter and distance
along the line are uniformly comparable.  Yet all these points lie
within $O(\Delta)$ of a union of at most
\[
 C M\Delta^{-x-\epsilon}
    \le C\Delta^{-b-x-\epsilon}
\]
grid cells, using the output lists separately for every sampled
\emph{record}.  Repeated records or repeated directions cause no
problem.  Apply Equation~\eqref{eq:RW} with its point exponent $b$
and line-family exponent $S$.  For any fixed $\omega>0$, choosing
$\epsilon$ small enough gives the lower covering bound
\[
 \Delta^{-b-\min\{1,S,(S+b)/2\}+\omega}.
\]
Comparison with the upper count excludes every strict deficit in
Equation~\eqref{names:score-projection-conclusion}.

For completeness, suppose mass bounds are supplied only at radii
$\Delta^{\sigma_j}$ on a grid with successive gaps at most $\gamma$,
including the two endpoints.  Enclose a ball of arbitrary radius by a
bounded number of cells at the next coarser tested radius.  The
additional loss is at most $\Delta^{-2\gamma}$, since $S\le2$ and
$b\le1$.  Choose $\gamma$ before the other small losses to fit the
strict deficit.  Finally, errors of size $\Delta^{1-\tau}$ are handled
by rerunning the proof at that effective final resolution.  The output
count exponent becomes $x/(1-\tau)$, and the small power constants are
rescaled by the same factor; the ball powers $S,b$ are unchanged.  The
conclusion converges to the displayed one as $\tau\downarrow0$.

For the last assertion, the restriction estimate in
Lemma~\ref{names:score-law-change} has upper bound
$I(Z;W\mid D)/c_\Delta+h(c_\Delta)/c_\Delta$.  This is
$o(\log(1/\Delta))$ under the stated assumptions, since
$h(c)/c\le\log(1/c)+1$.  Selecting parents with
$c_D\ge c_\Delta/2$ loses at most half the restricted probability and
adds only an exponent-zero normalization factor to the mass bounds.
The rest of the proof is unchanged.
\end{proof}

\begin{remark}[Obtaining the hypotheses from scores]
\label{names:score-projection-trimming}
If at a sample the source prefix score at radius $\Delta^\sigma$ is
at least $S\sigma-\epsilon$, retain only samples with this property.
Every occupied source cell in the retained law then has original
conditional mass at most $\Delta^{S\sigma-\epsilon}$; restriction can
only reduce its unnormalized mass.  Intersecting finitely many such
conditions gives the source bound in
Equation~\eqref{names:score-projection-masses}.  The same argument
applies to direction cells and to unnormalized direction bounds first
proved under finer side data and then integrated to $D$.

Conversely, a projection score at most $x+\epsilon$, conditionally on
the full record and $D$, places the projection in a list of at most
$\Delta^{-x-\epsilon}$ popular cells: each listed cell has probability
at least $\Delta^{x+\epsilon}$.  Thus projection-score inequalities
give precisely the lists required by
Proposition~\ref{names:projection-rule}.  Bin the proposed rates on a
positive-probability strict violation before making these finite
tests.  Corollary~\ref{names:score-future-prefix} supplies the needed
linear prefix rates also when the source is conditioned on a future
reading of another array.  No bound on the total number of source
cells is assumed.
\end{remark}

\subsection{Radial spreading and pin products}
\label{names:radial-section}

Projection tests whose directions come from two marks require information
about the lines joining those marks. The estimates below turn the time
marginal bounds into a dichotomy for tube masses, and provide a
product-coordinate estimate for the low-exponent alternative.
We begin with a finite-scale radial argument. It is useful to keep the
measure unnormalized: no assertion of conditional
independence will be made after restricting to a relation of pairs.
All measures in this subsection are supported in a fixed bounded
subset of $\R^2$.  Constants may depend on that subset.  We write
$\mathfrak l(x,y)$ for the unoriented line through distinct points
$x,y$, and $N_R(\ell)$ for its $R$-neighborhood.

\begin{lemma}[Finite radial exclusion]\label{names:radial-finite}
Assume the discretized Furstenberg estimate \eqref{eq:RW}.
For $0<A<d\leq1$ there are $\eta_0>0$ and
$\Delta_0>0$ with the following property.  There do not exist
$0<\Delta<\Delta_0$, a Borel measure $\mu$ of mass at most one,
and a symmetric measurable relation $E$ of distinct pairs such that
\begin{align}
 \mu(B(x,R))&\leq \Delta^{-\eta_0}R^d
       &&(\Delta\leq R\leq1),\label{names:radial-ball}\\
 (\mu\times\mu)(E)&\geq\Delta^{\eta_0},\label{names:radial-density}\\
 \mu(N_\Delta(\mathfrak l(x,y)))&\geq\Delta^{A+\eta_0}
       &&((x,y)\in E),\notag\\
 \mu(N_R(\mathfrak l(x,y)))&\leq\Delta^{-\eta_0}R^A
       &&((x,y)\in E,\ \Delta\leq R\leq1).
 \tag{rad-assume}\label{eq:rad-assume}
\end{align}
Replacing fixed constants in the radii or bounds by other fixed
constants does not change the conclusion, after reducing $\eta_0$
and $\Delta_0$.
\end{lemma}

\begin{proof}
Put
\[
 F(A,d)=A+\min\{1,d,(A+d)/2\},\qquad
 g=F(A,d)-2A>0.
\]
Choose $\kappa,\epsilon_0>0$ so small that
$\kappa+\epsilon_0<g/4$ and $\kappa<(1+A/d)/2$. We choose
a separation parameter $\lambda>0$ in the clustered case below.
We prove the assertion with $\eta_0=\eta$, where $\eta$ is chosen
last, sufficiently small relative to these parameters. The exact
requirements on $\eta$ are collected at the end of the proof.

Cover the bounded space of lines meeting the support by finitely
many slope--intercept charts and partition each chart into
$\Delta$-squares.  Assign each witness line to one such square
$c$, and let $T_c$ be a fixed $C\Delta$-tube containing the
$\Delta$-neighborhoods of every line assigned to $c$.  Define
\[
 J_c=\iint_{T_c\times T_c}
             \mathbf 1_{\{|u-v|\geq\Delta^\kappa\}}
             \,d\mu(u)\,d\mu(v).
\]
A pair at separation $r\geq\Delta^\kappa$ belongs to at most
$C(1+r^{-1})\leq C\Delta^{-\kappa}$ such pairs of footprints.
Indeed its allowable line directions occupy intervals of total
length $O(\Delta/r)$, while the intercept has only $O(1)$
choices at width $\Delta$ for each direction cell.  Consequently
\[
 \sum_cJ_c\leq C\Delta^{-\kappa}.
\]
Call a cell spread if $J_c\geq\Delta^{2A+\epsilon_0}$.
There are at most
\begin{equation}\label{names:radial-spread-count}
 C\Delta^{-2A-\epsilon_0-\kappa}
\end{equation}
spread cells.

For spread cells, this upper count will contradict the Furstenberg
lower count $\Delta^{-F(A,d)+g/4}$. In a nonspread cell, the small
mass of separated pairs instead forces most of the tube's mass
into one ball of radius $\Delta^\kappa$. We treat these two cases
in that order.

Suppose first that spread cells account for at least half of the
mass in \eqref{names:radial-density}.  Pass to one line chart,
losing a fixed factor, and discard pins whose partner mass in this
relation is less than $\Delta^{2\eta}$.  The retained pins have
mass at least $\Delta^{2\eta}$.  Their normalized distribution
has ball bounds $\Delta^{-3\eta}R^d$, with a further fixed
constant harmless.  For each retained pin $x$, push its restricted
partner measure to the line-parameter points
$\mathfrak l(x,y)$ and normalize.  This probability has ball
bounds
\[
 C\Delta^{-3\eta}R^A,\qquad \Delta\leq R\leq1.
\]
To check this, choose one witness line in a parameter ball which
meets the support of the pushforward.  Every partner belonging
to that ball lies in a $CR$-neighborhood of this witness line,
so \eqref{eq:rad-assume} applies.  Values $CR>1$ are covered by
the mass bound and a fixed constant.

Under planar duality the parameter points for a fixed pin lie on
its dual line.  In the chosen bounded chart the map from pins to
dual-line parameters is bi-Lipschitz.  Thus \eqref{eq:RW} applies
with point exponent $A$ and line exponent $d$, giving at least
$\Delta^{-F(A,d)+g/4}$ parameter cells.  This contradicts
\eqref{names:radial-spread-count}, since
$\epsilon_0+\kappa<g/4$ and $\Delta$ is sufficiently small.

It remains to treat witness pairs assigned to nonspread cells.
For any such occupied cell put $m_c=\mu(T_c)$; then
$m_c\geq\Delta^{A+\eta}$.  Averaging over $z\in T_c$ gives
a point $z_c\in T_c$ such that
\[
 \mu(T_c\setminus B(z_c,\Delta^\kappa))
       \leq J_c/m_c
       \leq\Delta^{A+\epsilon_0-\eta}.
\]
Hence the cluster
\[
 K_c=T_c\cap B(z_c,\Delta^\kappa)
\]
has mass at least $\tfrac12\Delta^{A+\eta}$, and therefore
at least $\Delta^{A+2\eta}$ when $\Delta$ is sufficiently small.

Choose
\[
 0<\lambda<\kappa/10,
 \qquad 10\lambda<\kappa(d-A).
\]
Discard witness pairs at distance less than $\Delta^\lambda$.
Their total mass is at most $\Delta^{d\lambda-\eta}$, which
is negligible compared with $\Delta^\eta$.  Since
$\lambda<\kappa$, at least one endpoint of every remaining pair
has distance at least $\tfrac14\Delta^\lambda$ from its
cluster.  Symmetry permits orienting a relation of mass at least
a fixed multiple of $\Delta^\eta$ so that its first endpoint
$x$ has this property.  Retain pins whose partner mass is at
least $\Delta^{2\eta}$; the set of retained pins again has
mass at least $\Delta^{2\eta}$.

For each retained pin take a maximal separated collection of its
witness directions, with angular separation
\[
 H=\Delta^{1-2\lambda}.
\]
A ball of directions of radius $CH$ accounts for partner mass
at most
$C\Delta^{-\eta}H^A\leq\Delta^{A(1-2\lambda)-2\eta}$,
by the tube upper bound centered at one witnessing line.
It follows that the collection has at least
\begin{equation}\label{names:radial-number-clusters}
 \Delta^{-A+2A\lambda+4\eta}
\end{equation}
members, with harmless rounding of this lower bound.  Assign to
each member the cluster of its witness cell.  These clusters are
disjoint: every point of such a cluster has distance
$\gtrsim\Delta^\lambda$ from $x$, so its direction from $x$
differs from the witness direction by $O(\Delta^{1-\lambda})$,
which is much smaller than $H$.

Let $E'$ consist of the new pairs $(x,y)$ with $y$ in one of
these chosen clusters.  The choices can be made measurably by
a fixed finite grid and a deterministic greedy ordering.  Using
\eqref{names:radial-number-clusters} and the cluster mass gives
\[
 (\mu\times\mu)(E')
    \geq\Delta^{2A\lambda+8\eta}\geq\Delta^{3\lambda}.
\]
Retain with each new pair its original witness line $\ell_0$
and cluster $K$.  We have
\begin{equation}\label{names:radial-cluster-transfer}
 K\subset B(y,C\Delta^\kappa)
              \cap N_{C\Delta}(\mathfrak l(x,y)),
 \qquad \mu(K)\geq\Delta^{A+2\eta}.
\end{equation}
Here $x\in\ell_0$, $y\in N_{C\Delta}(\ell_0)$, and
$|x-y|\gtrsim\Delta^\lambda$.  The angle of the two lines is
$O(\Delta^{1-\lambda})$.  On the cluster, whose diameter is
$O(\Delta^\kappa)$, rotation about $y$ costs at most
$O(\Delta^{1+\kappa-\lambda})=O(\Delta)$.  This proves
\eqref{names:radial-cluster-transfer} without thickening its
final tube by a power of $\Delta$.

Set $q=(1+A/d)/2$.  By \eqref{names:radial-ball},
\[
 \mu(B(y,\Delta^q))\leq\Delta^{(d+A)/2-\eta}
                         =o(\Delta^{A+2\eta}).
\]
Pigeonholing the dyadic annuli between radii $\Delta^q$ and
$C\Delta^\kappa$ therefore assigns to every pair of $E'$
a radius $D$ in this range for which its cluster contributes
at least $\Delta^{A+3\eta}$ mass in the annulus
$D\leq|z-y|\leq2D$.  There are $O(\log(1/\Delta))$
choices of $D$.  For one common $D$, the corresponding relation
has mass at least $\Delta^{4\lambda}$.  Choose $y$ so that
its set of pins has mass at least $\Delta^{5\lambda}$, and put
\[
 R=\Delta/D.
\]
This pin submeasure has angular ball bounds
\begin{equation}\label{names:radial-new-pin-bound}
 \mu\{x:(x,y)\in E',\;
             \angle(\mathfrak l(x,y),\ell)\leq CR\}
       \leq\Delta^{-2\eta}R^A
\end{equation}
whenever the angular ball meets the retained pin set; increasing
the exponent loss by $\eta$ absorbs any additional fixed constant.
In fact choose one pin in that ball and use its original witness
line $\ell_0$.  The new line differs from $\ell_0$ by
$O(\Delta^{1-\lambda})$, whereas
$R\gtrsim\Delta^{1-\kappa}$ and $\lambda<\kappa$.
Every pin in the angular ball is consequently in
$N_{C'R}(\ell_0)$.  The original upper bound in
\eqref{eq:rad-assume}, not an upper bound asserted for a newly
conditioned measure, proves \eqref{names:radial-new-pin-bound}.

A maximal collection of pin directions separated by a sufficiently
large constant times $R$ has at least
$\Delta^{5\lambda+3\eta}R^{-A}$ members.  Their
$C\Delta$-tubes in the annulus of radius $D$ about $y$ are
disjoint.  Each contains mass at least $\Delta^{A+3\eta}$
from its assigned cluster.  Consequently
\[
 \mu(B(y,2D))
   \geq \Delta^{5\lambda+6\eta}R^{-A}\Delta^A
   =\Delta^{5\lambda+6\eta}D^A.
\]
The ball upper bound, with its fixed-radius constant absorbed,
is at most $\Delta^{-2\eta}D^d$.  These inequalities imply
\[
 \Delta^{5\lambda+8\eta}
       \leq D^{d-A}
       \leq C\Delta^{\kappa(d-A)},
\]
contrary to the choice of $\lambda,\eta$ for small $\Delta$.

To meet all the losses used above, after choosing
$\kappa,\epsilon_0,\lambda$ in that order choose $\eta>0$ with
\[
 100\eta<\min\{\lambda,d\lambda,\epsilon_0,d-A\},
 \qquad 10\lambda+100\eta<\kappa(d-A),
\]
and smaller than the loss allowed by \eqref{eq:RW} with conclusion
loss $g/4$. These are finitely many compatible requirements.
Finally choose $\Delta_0$ small enough to absorb the fixed constants
and logarithmic factors in both cases.
\end{proof}

\begin{proposition}[Radial dichotomy]\label{names:radial-dichotomy}
Let $\xi_i=(\theta_i,r_i)$ and $\xi_j=(\theta_j,r_j)$
be conditionally independent, identically distributed planar
marks given $Z$.  Suppose their bounded-depth records obey the
preceding score calculus, and their time marginals have prefix
mass exponent at least $d>0$ given $Z$, with $d\leq1$.
More precisely, at every prescribed finite collection of depths,
discarding a set of vanishing average probability makes the
unnormalized conditional time-interval bounds hold with exponent
$d$ and arbitrarily small losses.  Assume \eqref{eq:RW}.

Trim at the required depths to measures $\mu_Z$ dominated by
the conditional mark laws.  Pass to a joint subsequential limit
of the needed scores and of
\[
 e(t)=-\log_{1/\rho}
        \mu_Z(N_{\rho^t}(\mathfrak l(\xi_i,\xi_j))),
 \qquad t>0\text{ rational}.
\]
Then, almost surely,
\begin{equation}
 e_*:=\liminf_{t\downarrow0}\frac{e(t)}{t}
          \in\{0\}\cup[d,\infty].
 \tag{rad}\label{eq:rad}
\end{equation}
If, in addition, the conditional tube supremum satisfies a
uniform affine-line avoidance estimate of some positive exponent
in these depth units, then $e_*\geq d$ almost surely.
\end{proposition}

\begin{proof}
Discard at each tested depth the time bins which violate the
asserted mass bound.  Their average mass tends to zero; finitely
many such trimmings have the same property.  Each time interval
is covered by a bounded number of comparable bins, so the trimmed
time marginal, and hence every planar ball, has the required
unnormalized bound.  Increasing finite grids can be chosen
diagonally.  Only grids fixed before the $\rho$-limit are used.

The displayed tube masses give legitimate limiting random
variables.  Indeed an endpoint ball of a sufficiently smaller
constant radius lies in the joining tube.  Among $O(\rho^{-2t})$
planar cells, the total trimmed mass of those having mass less
than $\rho^{2t+\alpha}$ is $O(\rho^\alpha)$.  The probability
that either endpoint was trimmed away tends to zero.  Thus the
exponents can be clipped on exceptional sets and jointly
subsequenced at countably many rational queries.  Nearby depth
queries handle constant changes of width.

If $0<e_*<d$ on a set of positive probability, choose an
essential value $A\in(0,d)$.  Given the tolerance in
Lemma~\ref{names:radial-finite}, restrict to
$|e_*-A|<\epsilon$, with $\epsilon$ sufficiently small.
At each such sample the lower inequalities
$e(u)\geq(A-2\epsilon)u$ hold for every sufficiently small
rational $u$, while
$e(t)\leq(A+2\epsilon)t$ holds at arbitrarily small rational
$t$.  A countable union therefore supplies a fixed small $t$
and a positive-probability event on which these inequalities
hold at $t$ and at every point of a prescribed finite grid
$\sigma t$, $0<\sigma\leq1$.

Choose that grid with mesh small relative to the tolerance of
Lemma~\ref{names:radial-finite}, including a first point so
small that the trivial mass bound fills the gap to radius one.
Monotonicity fills all intervening radii.  Transfer the strict
inequalities to the approximating laws and put $\Delta=\rho^t$.
The time bounds give \eqref{names:radial-ball}; the tube bounds
give \eqref{eq:rad-assume} with any prescribed sufficiently
small loss.  The event has fixed positive probability at this
stage.  Since both marks were sampled independently given $Z$,
its surviving pair mass is the average of the corresponding
$\mu_Z\times\mu_Z$ relation masses.  One can consequently
choose $Z$ with pair mass bounded below by a fixed positive
constant, in particular by $\Delta^{\eta_0}$ for small
$\Delta$.  The event and its transpose may be combined to make
the relation symmetric.  This contradicts
Lemma~\ref{names:radial-finite}.

A uniform positive tube-avoidance exponent gives $e_*>0$.
The dichotomy then gives the last assertion.
\end{proof}

The radial dichotomy bounds tubes under the original conditional mark
law. In the interior projection tests one also fixes the parent cell of
a moving mark and auxiliary angular data. The next lemma carries the
tube bounds to directions from that parent, keeping the source law and
the full pin record explicit.

\begin{lemma}[Parent-centered directions and comparison laws]
\label{names:parent-direction-transfer}
Let \(P\) be a law of a side variable \(Z\), finitely many full
records \(\boldsymbol V\), and a finite auxiliary name \(A\).
Suppose that the record marginal is conditionally iid given \(Z=z\),
with law \(\nu_z\). Each record has a bounded planar projection
\(\pi_z(V_i)\); write \(\mu_z=(\pi_z)_\#\nu_z\) for its law.
The full record may retain additional entries, including a stored
curvature slope. Let \(S\) be a function of \(Z\), and set
\[
 R(dZ,d\boldsymbol V,dA)
   =P(dZ,d\boldsymbol V)P(dA\mid S).
\]
Assume \(\operatorname{KL}(P\Vert R)=o(\log(1/\delta))\).
Choose distinct records \(V_x,V_y\), and write
\(x=\pi_z(V_x)\), \(y=\pi_z(V_y)\).
Let \(K\) be the depth-\(b\)
parent cell of \(y\), with representative \(c_K\), where \(b>0\).
For parents of positive mass put
\(\mu_z^K=\mu_z|_K/\mu_z(K)\).
The common data are
\[
 D=(Z,K,A,\text{the held full records other than }V_x,V_y).
\]
Fix \(t>0\) sufficiently small compared with \(b\), and a finite
grid of radii \(r=\delta^{\sigma t}\), \(0<\sigma\leq1\).
Let \(T_z\) be a measurable trimming set in the full record space
and put \(\mu_z^T=(\pi_z)_\#(\nu_z|_{T_z})\).
In particular the trimming may depend only on the planar projection.
Let \(W_y=\psi_D(\pi_z(V_y))\) be a finite Borel source name in the
tested prefix family. Under \(R\), its conditional law given \(D\)
is the image of \(\mu_z^K\) under \(\psi_D\); in applications it is a
bin of the rescaled child increment of the moving mark in \(K\).
On an event \(G\), suppose that \(V_x\in T_z\), that
\(\abs{x-y}\geq\delta^\chi\) with \(b-\chi>t\), and that
the original witness lines satisfy the tested bounds
\[
 \mu_z^T(N_{Cr}(\mathfrak l(x,y)))
       \leq C\delta^{-\eps t}r^{d-\eps}.
\]
Here the fixed enlargement constants needed below are included
among the finite tests. Then, with \(g_K(x)\) the direction from
\(c_K\) to \(x\),
\begin{equation}
 R(G,\ g_K(x)\in J(r)\mid D)
       \leq C\delta^{-\eps t}r^{d-\eps}
 \label{names:parent-direction-bound}
\end{equation}
for every tested direction interval \(J(r)\). A fixed swap or sign
change of its coordinate representation preserves the assertion.

After deleting a set of \(P\)-probability tending to zero, the
corresponding unnormalized conditional bound under \(P\) holds
with factor \(C\delta^{-3\eps t}\). Source prefix bounds for \(W_y\)
supplied by the original conditional
law \(\mu_z(\cdot\mid K)\), with their own source trimming when
needed, transfer in the same way. Moreover,
\[
 I_P(W_y;V_x\mid D)
       \leq\operatorname{KL}(P\Vert R)
       =o(\log(1/\delta^t)).
\]
All assertions fix \(b,t\), the finite grid, and the tested
positive-probability event before sending the information error to
zero. No lower bound for \(\mu_z(K)\) is required.
\end{lemma}

\begin{proof}
Since \(S\) is known from \(Z\), conditioning on \(A\) does not
change the full record kernel under \(R\). Conditional iid sampling
also allows the other full records to be held. On parents of positive
probability, the exact conditional pair law is consequently
\[
 R(dV_x,dV_y\mid D)=\nu_z(dV_x)\nu_z^K(dV_y),\qquad
 \nu_z^K=\frac{\nu_z|_{\pi_z^{-1}(K)}}{\mu_z(K)}.
\]
Its moving planar marginal is
\(\mu_z^K=(\pi_z)_\#\nu_z^K=\mu_z|_K/\mu_z(K)\).
This factorization holds before selecting \(G\). Take the jointly
Borel inner version of \(G\) from Lemma~\ref{cyl:borel-paths}
under the sum of the \(P\)- and \(R\)-laws on the full record tuple.
It preserves their masses and the pointwise witness bounds; all
conditional assertions here are for almost every common datum.
Its direction
marginal on \(G\) is
\begin{align*}
 &\int\mathbf1_{J(r)}(g_K(\pi_z(V_x)))\nu_z(dV_x)
          \int\mathbf1_G(D,V_x,V_y)\nu_z^K(dV_y)\\
 &\quad\leq\mu_z^T(\mathcal W_{D,J}),
\end{align*}
where \(\mathcal W_{D,J}\) is the set of planar projections of
surviving pin records in that direction interval having some
full-record witness \(V_y\) with projection in \(K\).
The inner probability is at most one. In particular the bound
does not contain a factor \(\mu_z(K)^{-1}\).

If \(\mathcal W_{D,J}\) is nonempty, choose one full-record
witness and denote its planar projections by \((x_0,y_0)\).
Since \(y_0\) is within \(O(\delta^b)\) of
\(c_K\) and the pair is separated, the directions of
\(\mathfrak l(x_0,y_0)\) and \(\mathfrak l(x_0,c_K)\)
differ by \(O(\delta^{b-\chi})=o(r)\).
Every participating pin lies in a fixed bounded region and makes
angle \(O(r)\) with that witness direction about \(c_K\).
It is therefore in \(N_{Cr}(\mathfrak l(x_0,y_0))\).
The original trimmed tube bound proves
Equation~\eqref{names:parent-direction-bound}. Notice that the
source law is still \(\mu_z^K\); it has not been replaced by
the conditional law of the trimmed measure.

For the change of law use the full finite tuple carrying the
event, the common data, and all tested source and pin names.
Its relative entropy is at most the stated budget by data
processing. Write \(f=dP/dR\) on this tuple and
\(f_D=dP_D/dR_D\). The likelihood estimates in
Lemma~\ref{names:score-law-change} show, for fixed
\(\eps,t>0\), that outside a set of vanishing \(P\)-probability,
\[
 f\leq\delta^{-\eps t},\qquad f_D\geq\delta^{\eps t}.
\]
The conditional likelihood on the retained event is at most
\(\delta^{-2\eps t}\). Integrating it over either the source
or the pin proves the claimed unnormalized marginal bounds
under \(P\). Intersecting further with a popular-output event
only decreases these masses.

Finally, decomposition of relative entropy conditional on \(D\),
followed by its decomposition against the product kernel
\(\nu_z\times\nu_z^K\), bounds
\(I_P(V_y;V_x\mid D)\) by
\(\operatorname{KL}(P\Vert R)\). The name \(W_y\) is a function of
\(V_y,D\), so data processing gives the asserted bound on
\(I_P(W_y;V_x\mid D)\), also with the full pin record
retained in the projection lists. The scale conversion is valid
because \(t\) is fixed. If a strict violation is retained on an
event of probability \(c(t)>0\), take the earlier information
error negligible compared with
\(c(t)t\log(1/\delta)\) before invoking the positive-event
restriction rule. This uses no uniform assertion as
\(t\downarrow0\).
\end{proof}

For the linear-fiber estimates, the projection direction is parametrized
by the scalar $(\theta_j-\theta_i)(r_j-r_i)$. We next control this
scalar on witness sets whose joining tubes have unusually large mass,
then combine that estimate with the high radial alternative.

\begin{lemma}[A low-exponent pin estimate]\label{names:radial-low-product}
Fix $0<d\leq1$.  There is a constant $C_d$ with the following
property, for every fixed $\epsilon>0$ and all sufficiently
small $\Delta>0$ depending on $\epsilon$.  Let $\mu$ have
mass at most one and time marginal
\[
 \mu\{(\theta,r):\theta\in I\}
           \leq\Delta^{-\epsilon}|I|^d
 \quad(\Delta\leq|I|\leq1).
\]
Fix a pin $x=(\theta_x,r_x)$ in the same fixed bounded set,
and let $W_x$ be any subrelation
of points $y=(\theta_y,r_y)$ for which
\[
 |\theta_y-\theta_x|,\ |r_y-r_x|\geq\Delta^\epsilon,
 \qquad
 \mu(N_\Delta(\mathfrak l(x,y)))\geq\Delta^\epsilon.
\]
For $f_x(y)=(\theta_y-\theta_x)(r_y-r_x)$, every interval
$I$ of length $\Delta^\sigma$, $0\leq\sigma\leq1$, satisfies
\begin{equation}\label{names:radial-low-product-bound}
 \mu\{y\in W_x:f_x(y)\in I\}
            \leq C_d\Delta^{d\sigma-C_d\epsilon}.
\end{equation}
Finite grids of prefix hypotheses suffice, with their mesh
added to the exponent loss.
\end{lemma}

\begin{proof}
It suffices to treat $3\epsilon/d<1$: by increasing $C_d$,
the conclusion is immediate from the total mass bound for
larger $\epsilon$.
The time bound also bounds planar balls.  Put
$a=\Delta^{3\epsilon/d}$.  A ball of radius $a$ about $x$
has mass at most $C\Delta^{2\epsilon}$, so at least half
the mass of each witness tube is outside that ball, for small
$\Delta$.  Choose a maximal family of witness directions
separated by $C\Delta/a$.  Outside $B(x,a)$ their
$\Delta$-tubes are disjoint when $C$ is sufficiently large.
The number of chosen directions is at most
$C\Delta^{-\epsilon}$.  Maximality places all witness points
within $O(\Delta/a)$ of one of the representative lines.

Each representative is a witness, so its slope $m$ satisfies
$c\Delta^\epsilon\leq|m|\leq C\Delta^{-\epsilon}$.
On its enlarged tube,
\[
 f_x(y)=m(\theta_y-\theta_x)^2
                   +O(\Delta^{1-C_d'\epsilon}).
\]
For the points under consideration,
$|\theta_y-\theta_x|\geq\Delta^\epsilon$.
On each of the two possible signs of this difference, the
inverse of the quadratic has derivative at most
$C\Delta^{-2\epsilon}$.  Thus $f_x(y)\in I$ restricts
$\theta_y$ to at most two intervals of total length
$C\Delta^{\sigma-C_d''\epsilon}$.  Apply the time-marginal
bound and sum over at most $C\Delta^{-\epsilon}$
representative lines.  This gives
\eqref{names:radial-low-product-bound}; intervals longer than
one are handled by the total mass bound.  Rounding the queried
lengths to a finite grid merely adds its mesh to the loss.
\end{proof}

\begin{proposition}[Pin-product alternatives]\label{names:radial-pin-product}
In addition to the hypotheses of
Proposition~\ref{names:radial-dichotomy}, assume the following.
The score calculus and the preceding conditional projection
rule apply to all fixed finite queries used below.  At almost
every parent depth $p$, the joint mark increment, conditional
on $Z$ and its own parent $\xi_{j,p}$, has prefix mass exponent
at least $d$ on every fixed finite grid of sufficiently small
layers $p\to p+t$.  Finally, for every $\chi>0$,
\[
 \Prob\bigl\{|\theta_i-\theta_j|\geq\rho^\chi,
                 \ |r_i-r_j|\geq\rho^\chi\bigr\}
             \longrightarrow1.
\]
Put $f=(\theta_j-\theta_i)(r_j-r_i)$ and
$\mathcal G_+=\{e_*\geq d\}$, $\mathcal G_0=\{e_*=0\}$.
Then the limiting cumulative scores satisfy
\begin{equation}\label{names:radial-high-product}
 i(f_p\mid Z,\xi_i)\geq dp
       \quad\hbox{on }\mathcal G_+
\end{equation}
at the tested depths.  Here $\xi_i$ denotes a record stored to
a sufficiently large fixed depth, chosen before the $\rho$-limit.

On $\mathcal G_0$ the following local substitute holds.  Any
strict violation of a differentiable layer rate on a set of
positive probability can be tested at a fixed small layer $t$
on a positive-probability subset where
$e(t)<\epsilon t$.  With $\Delta=\rho^t$, the unnormalized
witness sublaw of $f$, given $Z,\xi_i$, then has prefix
interval bounds of exponent $d$ with loss $O_d(\epsilon)$,
as in \eqref{names:radial-low-product-bound}.  The assertion is
for a fixed selected layer before earlier approximation and
information errors are sent to zero.
\end{proposition}

\begin{proof}
We first record precisely how a high radial bound produces a
direction bound.  On any finite grid of sufficiently small
depths $u=\sigma t$, a high witness satisfies
\[
 \mu_Z(N_{\rho^u}(\mathfrak l(\xi_i,\xi_j)))
                  \leq\rho^{(d-\epsilon)u},
\]
with a further arbitrarily small absolute loss when transferring
to the approximating laws.  For fixed $Z,\xi_i$, an angular bin
which meets the witness set has all its witnessing partners
inside a constant enlargement of one witness tube.  This gives
the same bound for the unnormalized partner-direction sublaw.
One may instead fix $Z$ and a $j$-parent bin at depth $p\gg t$.
Use its representative as the center.  After discarding pairs
closer than a sufficiently small fixed power, the change of
center is negligible at every queried angular width.  An
angular bin with a witness again places all its participating
pins inside an enlargement of that witness tube.  Thus the
pin-direction sublaw, unnormalized and with this parent held,
has the same exponent $d$.  These are bounds on a restricted
measure, not claims that the two marks remain independent
after the witness restriction.

The cumulative score of $f$ given $Z,\xi_i$ is Lipschitz,
starts at zero, and is differentiable at almost every depth.
If \eqref{names:radial-high-product} fails on a set of positive
probability, Fubini's theorem and the fundamental theorem of
calculus provide a fixed differentiability depth $p>0$ and a
positive-probability high set where the incremental rate is
strictly less than $d$.  Choose $t\ll p$ small enough to test
this strict deficit, together with a finite prefix grid on
which the preceding high tube bounds hold.

Adding $\xi_{j,p}$ to the conditioning can only lower the
tested limiting incremental score, by the score calculus.
The coarse product is known to bounded lists from this parent
and the pin.  After translating and dilating the $j$-parent,
the product increment is, up to $O(\rho^p)$ in dilated units,
the linear form with coefficient vector
\[
 (r_{j,\mathrm{ctr}}-r_i,
                  \theta_{j,\mathrm{ctr}}-\theta_i).
\]
Its norm is bounded below by an arbitrarily small power with
probability tending to one.  Since $t\ll p$, both its
quadratic remainder and the finite-record errors fit the
tested resolution, with arbitrarily small exponent loss.
The source increment has prefix exponent at least $d$ by
hypothesis.  The gradient directions, given $Z,\xi_{j,p}$,
have unnormalized exponent at least $d$ on the high relation:
they are the joining directions just estimated, with the two
coordinates interchanged.  Before restriction, the pin is
independent of the $j$-increment given this parent, since the
two marks were independent given $Z$.  The conditional
projection rule applies and gives projected rate at least
$\min\{d,(d+d)/2,1\}=d$, contradicting the strict deficit.
Integrating the almost-everywhere rate bound proves
\eqref{names:radial-high-product}.  Popular-cell trimming
converts it to the corresponding unnormalized interval-mass
bounds whenever needed.

On $\mathcal G_0$, the definition of a pointwise liminf supplies
arbitrarily small rational $t$ with $e(t)<\epsilon t$ at
each sample.  A strict violation of a differentiable limiting
rate persists at every sufficiently small layer at that sample.
A countable union therefore selects one fixed such $t$ on a
positive-probability violating set.  Impose the two scalar
separations at threshold $\rho^{\epsilon t}$; their exceptional
probability tends to zero in the earlier limits.  The trimmed
time marginal has the required bounds at the finite queried
depths.  At $\Delta=\rho^t$, the condition on $e(t)$ gives
the tube lower bound of
Lemma~\ref{names:radial-low-product}, with an arbitrarily
small additional loss.  That lemma proves the asserted
interval bounds at all the necessary prefixes.

Finally, these conditional bounds can be averaged over $Z$
and independently sampled auxiliary records when the parent
conditioning keeps the other pin fixed and permits that
mixture.  A different conditioning requires its own
justification.  In particular the selected low-set event
may have a small positive probability depending on $t$.
One first fixes this $t$ and its finite queries, and only
then makes the earlier information errors negligible relative
to that probability.  Nothing here permits transferring an
information budget through a restriction whose probability
is merely a power of the fine scale.
\end{proof}
